% !TEX program = xelatex
% =====================================================================
%  BÁO CÁO ĐỒ ÁN TRÍ TUỆ NHÂN TẠO, HỌC PHẦN AIP202
%  Khoa Công nghệ Thông tin, Trường Đại học Kiến trúc Đà Nẵng
%
%  Build:  latexmk -xelatex do-an.tex
%  Overleaf: Menu > Compiler > XeLaTeX, rồi bấm Recompile.
%
%  Khung tài liệu chuyển thể từ mẫu đồ án của Trường Đại học Nha Trang
%  (github.com/nd-hung/thesis-template, tác giả TS. Nguyễn Đình Hưng).
%  Toàn bộ nhận diện, phông chữ, bảng màu, trang bìa và cấu trúc chương
%  được thay theo chuẩn của DAU và theo đề cương học phần AIP202.
% =====================================================================
% \documentclass phải nằm ở file này, không được chuyển vào settings.tex,
% vì Overleaf dò file chính bằng cách tìm file có chứa \documentclass.
\documentclass[12pt, a4paper, oneside]{book}
% !TEX root = ../do-an.tex
% =====================================================================
%  styles/settings.tex
%  Các thiết lập chung, nạp từ file chính do-an.tex.
%  SINH VIÊN KHÔNG SỬA FILE NÀY.
%
%  Lưu ý: lệnh khai báo lớp tài liệu nằm ở file do-an.tex chứ không nằm
%  ở đây. Lý do: Overleaf dò file chính bằng cách tìm chuỗi khai báo lớp
%  tài liệu trong các file. Nếu chuỗi đó xuất hiện ở file này, Overleaf
%  sẽ biên dịch nhầm file này và báo lỗi "no legal end found", vì file
%  này chỉ chứa thiết lập chứ không chứa phần thân tài liệu.
% =====================================================================

% Toàn bộ nhận diện DAU: phông, màu, bìa, hộp trình bày, header
\usepackage{styles/dau}
\dauLayout

% Giãn dòng và khoảng cách đoạn theo quy cách đồ án
\linespread{1.25}
\usepackage[skip=8pt plus 1pt, indent=15pt]{parskip}

% Ghi chú: mẫu gốc dùng gói atbegshi để vứt trang đầu, vì trang bìa của
% mẫu gốc dựng bằng minipage nên sinh ra một trang trắng. Trang bìa ở
% đây dựng bằng môi trường titlepage nên không có trang trắng, và nếu
% giữ thủ thuật đó thì chính trang bìa sẽ bị vứt mất.


% =====================================================================
%  PHẦN EM ĐIỀN: thông tin đồ án
% =====================================================================
% Loại báo cáo
\crname{BÁO CÁO ĐỒ ÁN TRÍ TUỆ NHÂN TẠO}
% Học phần
\cmonhoc{Học phần AIP202 \textperiodcentered{} 2 tín chỉ}
% Tên đề tài. Dùng \\ để ngắt dòng cho cân đối trên bìa.
\ctname{XÂY DỰNG HỆ THỐNG GỢI Ý PHỐI ĐỒ}
% Họ tên sinh viên và mã số sinh viên
\cstuname{Giản Trần Phước Anh}
\cmssv{2351220045}
\cstunametwo{Lê Bạch Tấn Khang}
\cmssvtwo{2351220010}
% Lớp
\clopsv{CNTT 22A}
% Ngành
\csCouncil{Công nghệ thông tin}
% Giảng viên hướng dẫn
\csSupervise{TS. Đỗ Phúc Hảo}
% Năm học
\cttime{2025--2026}
% Khoa quản lý
\csdeptname{KHOA CÔNG NGHỆ THÔNG TIN}
% Nơi thực hiện
\csplace{Đà Nẵng}

\begin{document}

% Bìa trang trọng có khung kép. Muốn dùng bìa hiện đại theo bộ giáo
% trình của khoa thì thay dòng dưới bằng \coverpagehiendai
\coverpage

% ---------------------------------------------------------------------
%  PHẦN ĐẦU: đánh số trang bằng số La Mã thường
% ---------------------------------------------------------------------
\frontmatter
\dauPageStyle

% ---------------------------------------------------------------------
%  LỜI CAM ĐOAN
% ---------------------------------------------------------------------
\begin{declaration}

Chúng tôi xin cam đoan đồ án này là công trình do chính chúng tôi thực hiện dưới sự
hướng dẫn của \dauGVHD. Các nội dung nghiên cứu, số liệu và kết quả trình
bày trong đồ án là trung thực và chưa từng được công bố trong bất kỳ công
trình nào khác.

Mọi số liệu, hình ảnh, mã nguồn và kết quả của tác giả khác được sử dụng
trong đồ án đều được trích dẫn nguồn gốc rõ ràng trong phần tài liệu tham
khảo. Toàn bộ phần mềm và thư viện được dùng đều có giấy phép cho phép sử
dụng. Những phần có sử dụng công cụ trí tuệ nhân tạo hỗ trợ đã được nêu rõ
trong báo cáo.

Nếu phát hiện có bất kỳ sự gian lận nào, chúng tôi xin chịu hoàn toàn trách nhiệm
trước Khoa và Nhà trường.

\vspace{1.4em}
\noindent\okyten{\dauSinhVien}{Sinh viên thực hiện}\hfill
\okyten{\dauSinhVienHai}{Sinh viên thực hiện}

\end{declaration}

\begin{acknowledgments}

Chúng em xin gửi lời cảm ơn chân thành đến \dauGVHD, người đã tận tình định
hướng đề tài, góp ý và theo sát tiến độ của chúng em trong suốt quá trình
thực hiện đồ án. Chúng em cũng xin cảm ơn quý thầy cô Khoa Công nghệ Thông
tin, Trường Đại học Kiến trúc Đà Nẵng đã trang bị cho chúng em nền tảng kiến
thức để hoàn thành công việc này. Cuối cùng, chúng em xin cảm ơn gia đình và
bạn bè đã luôn động viên nhóm trong thời gian học tập.

\vspace{1.4em}
\noindent\okyten{\dauSinhVien}{Sinh viên thực hiện}\hfill
\okyten{\dauSinhVienHai}{Sinh viên thực hiện}

\end{acknowledgments}

% ---------------------------------------------------------------------
%  TÓM TẮT
% ---------------------------------------------------------------------
\begin{abstract}

Việc lựa chọn trang phục phù hợp thường đòi hỏi người dùng cân nhắc đồng
thời loại quần áo, màu sắc, phong cách và hoàn cảnh sử dụng. Đồ án tập trung
xây dựng hệ thống gợi ý phối đồ trên thiết bị di động nhằm hỗ trợ người dùng
nhận diện trang phục từ ảnh, quản lý tủ đồ cá nhân và lựa chọn các món đồ có
khả năng kết hợp với nhau. Hệ thống được định hướng gồm ứng dụng Flutter,
dịch vụ xử lý trí tuệ nhân tạo thông qua API và Firebase để lưu thông tin
người dùng, quần áo cùng các kết quả yêu thích.

Đầu vào của chức năng nhận diện là ảnh chứa một trang phục chính; đầu ra là
loại quần áo dự đoán và độ tin cậy. Trên cơ sở loại trang phục, màu sắc,
phong cách và dữ liệu trong tủ đồ, hệ thống dự kiến trả về tối đa ba phương
án phối đồ, đồng thời hạn chế kết quả trùng lặp. Đồ án cũng giới hạn chức
năng thử trang phục ảo ở kỹ thuật ghép ảnh hai chiều: ảnh quần áo có nền
trong suốt được đặt lên ảnh người dùng và có thể điều chỉnh vị trí, kích
thước, góc xoay. Cách tiếp cận này phù hợp với phạm vi học phần và không
được xem là mô hình sinh ảnh hay mô phỏng vật lý của vải.

Thành phần nhận diện sử dụng mô hình phát hiện đối tượng YOLO26n và giữ cố
định 21 lớp quần áo. Dữ liệu huấn luyện kết hợp 11.926 ảnh xuất từ Roboflow
với 52.712 ảnh DeepFashion In-shop; trong đó 52.711 ảnh DeepFashion được
chuyển thành công sang định dạng YOLO. Tập dữ liệu chuẩn bị cho phiên bản V3
gồm 59.463 ảnh huấn luyện, 6.902 ảnh xác thực và 6.620 ảnh kiểm thử. Ở lần
fine-tune cuối, mô hình đạt precision 55,18\%, recall 49,19\%, mAP@0,5 là
49,53\% và mAP@0,5:0,95 là 36,10\% trên tập xác thực. Thời gian phản hồi API
trên thiết bị Android thật chưa nằm trong phạm vi số liệu của lần chạy tích
hợp được sử dụng trong báo cáo.

\vspace{0.8em}
\noindent\textbf{Từ khoá:} gợi ý phối đồ, nhận diện quần áo, thị giác máy
tính, Flutter, thử trang phục ảo.

\end{abstract}


\tableofcontents
\listoffigures
\listoftables
% Bỏ dòng dưới nếu đồ án của em không có giải thuật nào
\listofalgorithms

% ---------------------------------------------------------------------
%  THÂN BÁO CÁO: sáu chương theo quy cách của học phần
% ---------------------------------------------------------------------
\mainmatter
\dauPageStyle

% ---------------------------------------------------------------------
%  CHƯƠNG 1 — MỞ ĐẦU
% ---------------------------------------------------------------------
\chapter{Mở đầu}
\label{ch:modau}

\begin{muctieu}
Chương mở đầu trình bày lý do lựa chọn đề tài, mục tiêu cần đạt, đối tượng,
phạm vi thực hiện và bố cục của báo cáo. Qua đó, người đọc có thể nắm được
bức tranh tổng quan về ứng dụng trí tuệ nhân tạo trong nhận diện quần áo,
gợi ý phối đồ và thử trang phục ảo.
\end{muctieu}

\section{Lý do chọn đề tài}

Trong đời sống hiện nay, nhu cầu lựa chọn trang phục phù hợp với hoàn cảnh,
màu sắc và phong cách cá nhân ngày càng được quan tâm. Tuy nhiên, nhiều người
gặp khó khăn trong việc xác định loại quần áo, lựa chọn những món đồ phù hợp
với nhau hoặc hình dung trang phục sẽ trông như thế nào khi mặc. Việc thử
trực tiếp nhiều bộ trang phục cũng tốn thời gian, đặc biệt khi mua sắm trực
tuyến hoặc quản lý số lượng lớn quần áo cá nhân.

Sự phát triển của trí tuệ nhân tạo, đặc biệt là thị giác máy tính
\en{(computer vision)} và học sâu \en{(deep learning)}, cho phép máy tính
phân tích hình ảnh và nhận diện các đối tượng trong ảnh với độ chính xác
ngày càng cao. Những công nghệ này có thể được ứng dụng để nhận diện loại
quần áo, trích xuất một số đặc điểm như màu sắc và kiểu dáng, từ đó hỗ trợ
người dùng lựa chọn trang phục phù hợp.

Bên cạnh đó, điện thoại thông minh đã trở thành thiết bị phổ biến và thuận
tiện để chụp ảnh, lưu trữ thông tin cũng như sử dụng các dịch vụ thông minh.
Vì vậy, việc xây dựng một ứng dụng di động tích hợp trí tuệ nhân tạo có thể
giúp người dùng nhận diện quần áo, quản lý tủ đồ, nhận gợi ý phối đồ và thử
trang phục ảo ngay trên điện thoại.

Từ những nhu cầu thực tế trên, đề tài ``Xây dựng hệ thống gợi ý phối đồ''
được lựa chọn. Trong phạm vi đồ án, hệ thống được định hướng dưới dạng ứng
dụng di động Flutter, kết hợp nhận diện quần áo, quản lý tủ đồ, gợi ý phối
đồ và thử trang phục ảo ở mức cơ bản. Đề tài không chỉ có ý nghĩa thực tiễn
mà còn tạo điều kiện vận dụng các kiến thức về trí tuệ nhân tạo, xử lý ảnh,
phát triển ứng dụng di động, xây dựng API và quản lý dữ liệu.

\section{Mục tiêu của đồ án}

Mục tiêu tổng quát của đồ án là xây dựng một ứng dụng di động có khả năng
nhận diện quần áo từ hình ảnh, đề xuất cách phối đồ và hỗ trợ người dùng thử
trang phục ảo ở mức cơ bản.

Các mục tiêu cụ thể của đồ án gồm:

\begin{itemize}[leftmargin=1.6em, itemsep=3pt]
  \item Xây dựng mô hình trí tuệ nhân tạo có khả năng phân loại một số nhóm
        quần áo phổ biến từ ảnh do người dùng chụp hoặc tải lên.

  \item Đánh giá và so sánh ít nhất hai mô hình học sâu phù hợp với bài toán
        nhận diện hình ảnh, từ đó lựa chọn mô hình có sự cân bằng giữa độ
        chính xác và tốc độ xử lý.

  \item Hướng tới độ chính xác nhận diện từ 80\% trở lên trên tập kiểm thử
        và thời gian phản hồi không quá 3 giây trong điều kiện kiểm thử ổn
        định.

  \item Xây dựng chức năng gợi ý tối đa ba cách phối đồ phù hợp dựa trên loại
        trang phục, màu sắc và phong cách, đồng thời hạn chế các kết quả bị
        trùng lặp.

  \item Xây dựng chức năng thử trang phục ảo ở mức cơ bản bằng cách đặt ảnh
        quần áo có nền trong suốt lên ảnh người dùng và cho phép điều chỉnh
        vị trí, kích thước và góc xoay.

  \item Phát triển ứng dụng di động bằng Flutter, kết nối với hệ thống xử lý
        trí tuệ nhân tạo thông qua API.

  \item Xây dựng các chức năng hỗ trợ như đăng ký, đăng nhập, quản lý tủ đồ
        cá nhân, lưu trang phục yêu thích và lưu kết quả phối đồ.
\end{itemize}

\section{Phạm vi và đối tượng}

\subsection{Đối tượng nghiên cứu}

Đối tượng nghiên cứu của đồ án gồm:

\begin{itemize}[leftmargin=1.6em, itemsep=3pt]
  \item Các hình ảnh chứa một trang phục chính như áo thun, áo sơ mi, áo
        khoác, quần dài, quần ngắn, váy hoặc đầm.

  \item Các mô hình học sâu có khả năng phân loại hình ảnh quần áo và phù
        hợp với hệ thống có tài nguyên tính toán giới hạn.

  \item Các nguyên tắc phối hợp trang phục dựa trên loại quần áo, màu sắc
        và phong cách cơ bản.

  \item Phương pháp xác định vị trí cơ thể để đặt lớp ảnh trang phục lên ảnh
        người dùng.

  \item Quy trình phát triển ứng dụng di động bằng Flutter và kết nối ứng
        dụng với dịch vụ xử lý trí tuệ nhân tạo.
\end{itemize}

\subsection{Phạm vi thực hiện}

Trong phạm vi đồ án, hệ thống tập trung vào các chức năng sau:

\begin{itemize}[leftmargin=1.6em, itemsep=3pt]
  \item Nhận ảnh quần áo từ camera hoặc thư viện ảnh của điện thoại.

  \item Nhận diện một số loại quần áo phổ biến đã được xác định trong bộ dữ
        liệu huấn luyện.

  \item Trả về tên loại quần áo dự đoán và độ tin cậy của kết quả.

  \item Gợi ý ba phương án phối đồ dựa trên các thuộc tính cơ bản của trang
        phục và dữ liệu trong tủ đồ cá nhân.

  \item Hỗ trợ thử trang phục ảo bằng kỹ thuật ghép ảnh hai chiều
        \en{(2D overlay)}, không sử dụng mô hình sinh ảnh.

  \item Cho phép người dùng điều chỉnh thủ công vị trí, kích thước và góc
        xoay của trang phục khi thử đồ.

  \item Lưu thông tin người dùng, hình ảnh quần áo, tủ đồ và các kết quả yêu
        thích bằng Firebase.
\end{itemize}

Đồ án chưa tập trung vào việc nhận diện đồng thời nhiều loại trang phục
trong một ảnh phức tạp, xác định thương hiệu, chất liệu hoặc kích cỡ chính
xác của quần áo. Chức năng thử trang phục ảo chỉ mô phỏng bằng lớp ảnh hai
chiều nên chưa thể tái tạo chính xác độ rủ, nếp gấp, ánh sáng hoặc sự thay đổi
hình dạng của trang phục theo cơ thể người dùng.

Ứng dụng được xây dựng chủ yếu để thử nghiệm trên thiết bị Android. Hệ thống
chưa hướng tới triển khai thương mại, chưa tích hợp chức năng thanh toán và
chưa thay thế hoàn toàn việc thử trang phục trực tiếp.

\section{Bố cục báo cáo}

Phần còn lại của báo cáo được tổ chức như sau. Chương~\ref{ch:coso} trình
bày cơ sở lý thuyết về trí tuệ nhân tạo, thị giác máy tính, mô hình nhận diện
ảnh và các công nghệ được sử dụng. Chương~\ref{ch:thietke} phân tích yêu cầu
và thiết kế kiến trúc, dữ liệu, chức năng của hệ thống. Chương~\ref{ch:caidat}
mô tả quá trình xây dựng mô hình AI, phát triển ứng dụng Flutter và tích hợp
các thành phần. Chương~\ref{ch:ketqua} trình bày kết quả thực nghiệm, kiểm
thử và đánh giá hệ thống. Cuối cùng, Chương~\ref{ch:ketluan} tổng kết những
kết quả đã đạt được, nêu các hạn chế và đề xuất hướng phát triển trong tương
lai.

% ---------------------------------------------------------------------
%  CHƯƠNG 2 — CƠ SỞ LÝ THUYẾT VÀ CÔNG NGHỆ
% ---------------------------------------------------------------------
\chapter{Cơ sở lý thuyết và công nghệ}
\label{ch:coso}

\begin{muctieu}
Chương này trình bày nền tảng của bài toán nhận diện quần áo, gợi ý phối đồ,
thử trang phục ảo và các nhóm công nghệ được xác định trong phạm vi đồ án.
\end{muctieu}

\section{Kiến thức nền tảng}

\begin{dinhnghia}[Phát hiện trang phục]
Phát hiện trang phục là bài toán xác định hộp giới hạn, nhãn lớp và độ tin
cậy của một hoặc nhiều món đồ xuất hiện trong ảnh. Với mỗi dự đoán, mô hình
trả về hộp $B_p$, nhãn $y \in \mathcal{Y}$ và một điểm tin cậy; trong đồ án,
$\mathcal{Y}$ gồm 21 lớp quần áo và phụ kiện.
\end{dinhnghia}

YOLO thực hiện phát hiện đối tượng trong một lần suy luận, kết hợp định vị và
nhận dạng lớp trong cùng mạng nơ-ron~\cite{Redmon2016}. Notebook của đồ án
khởi tạo từ checkpoint \code{yolo26n.pt}~\cite{UltralyticsYOLO26}, sau đó fine-tune tuần tự qua các
phiên bản V1, V2 và V3 bằng thư viện Ultralytics. Ảnh đầu vào được đưa về
kích thước $640\times640$ pixel; cấu hình dữ liệu giữ nguyên thứ tự 21 lớp để
checkpoint và API diễn giải nhãn nhất quán.

Mức chồng lấp giữa hộp dự đoán $B_p$ và hộp thật $B_g$ được đo bằng IoU ở
Công thức~\eqref{eq:iou-report}:
\begin{equation}
\mathrm{IoU}(B_p,B_g)=\frac{|B_p\cap B_g|}{|B_p\cup B_g|}.
\label{eq:iou-report}
\end{equation}
IoU được dùng khi ghép dự đoán với nhãn thật và khi tính precision, recall,
mAP@0,5 và mAP@0,5:0,95. Các đại lượng này phù hợp với bài toán phát hiện hơn
accuracy của bài toán phân loại ảnh~\cite{COCO2014}.

DeepFashion In-shop cung cấp 52.712 ảnh cùng hộp giới hạn và thông tin sản
phẩm~\cite{DeepFashion2016}. Notebook V3 chuyển hộp giới hạn sang định dạng
YOLO, chia dữ liệu theo \code{item\_id} với hạt giống 42 để các góc chụp của
cùng một sản phẩm không rơi vào nhiều tập, rồi trộn với bộ dữ liệu Roboflow
ban đầu. Phép tăng cường trong cấu hình huấn luyện gồm biến đổi HSV, tịnh tiến,
co giãn, lật ngang và mosaic.

Gợi ý phối đồ khác với tìm kiếm sản phẩm tương tự vì hệ thống phải đánh giá
khả năng bổ sung giữa các món đồ thuộc nhiều nhóm khác nhau~\cite{McAuley2015}. Han và cộng sự
mô hình hoá một bộ trang phục như một chuỗi và học đồng thời quan hệ tương
thích cùng biểu diễn thị giác--ngữ nghĩa~\cite{Han2017}. Trong mã nguồn hiện
tại, luồng gợi ý là một chức năng không dùng ảnh và không nhận kết quả YOLO.
Backend gửi hồ sơ cơ thể, giới tính tự khai, thời tiết hoặc mùa/nhiệt độ,
phong cách và hoàn cảnh đến Gemini; kết quả có một nhận định, lý do và từ một
đến ba điều chỉnh. Luồng hỏi đáp dùng cùng bối cảnh và tối đa sáu lượt gần
nhất khi gọi mô hình.

Thử trang phục ảo có thể được thực hiện bằng mô hình sinh ảnh hoặc bằng ghép
ảnh hai chiều. VITON là một phương pháp học sâu chuyển một món đồ lên ảnh
người theo chiến lược từ thô đến tinh mà không sử dụng mô hình 3D
~\cite{VITON2018}. Ứng dụng hiện dùng MediaPipe Pose Landmarker ở backend để
trả landmark và các điểm neo cơ thể qua WebSocket. Flutter làm mượt tư thế,
đặt ảnh trang phục nền trong suốt theo nhóm cơ thể và vẫn cho phép người dùng
dịch chuyển, co giãn, xoay thủ công. Đây là ghép lớp theo tư thế, không phải
mô hình sinh ảnh nên không mô phỏng chính xác độ rủ, nếp gấp hay ánh sáng.

Ba phiên bản thực nghiệm dùng cùng họ YOLO thay vì ba kiến trúc độc lập. V1
khởi tạo trực tiếp từ \code{yolo26n.pt}; V2 fine-tune checkpoint V1 trên dữ
liệu được bổ sung một phần từ DeepFashion; V3 sử dụng toàn bộ 52.712 ảnh
DeepFashion và lặp lại ảnh train cũ để giảm quên các lớp phụ kiện. V2 và V3
dùng AdamW, batch 16, kích thước ảnh 640 và weight decay 0,0005. Vì vậy, mục
tiêu so sánh hai kiến trúc học sâu khác nhau chưa được đáp ứng bởi các tệp
huấn luyện hiện có.

\section{Công cụ và thư viện}

Các thành phần thực tế được đối chiếu từ \code{.metadata},
\code{pubspec.lock}, môi trường Python của backend và cấu hình huấn luyện,
sau đó tổng hợp ở Bảng~\ref{tab:congcu}. Phiên bản dùng khi huấn luyện trên
Colab được phân biệt với phiên bản đang cài trong backend triển khai.

\begin{table}[H]
\centering
\caption{Công cụ và thư viện sử dụng trong đồ án.}
\label{tab:congcu}
\footnotesize
\setlength{\tabcolsep}{4pt}
\begin{tabularx}{\textwidth}{>{\bfseries\color{daunavy}\raggedright\arraybackslash}p{2.8cm} >{\color{daumaroon}\raggedright\arraybackslash}p{3.4cm} >{\raggedright\arraybackslash}X}
\toprule
\rowcolor{gray100}\bfseries\color{daunavy}Thành phần & \bfseries\color{daumaroon}Phiên bản & \bfseries\color{daunavy}Vai trò \\
\midrule
Flutter / Dart & 3.47.1 stable / 3.13.1 & Xây dựng giao diện Android, điều hướng bốn tab Phân tích, Gợi ý, Hỏi đáp và Cá nhân \\
Firebase & Core 4.14.0; Auth 6.6.1; Firestore 6.9.0 & Khởi tạo ứng dụng, xác thực email/mật khẩu, lưu hồ sơ tài khoản và hồ sơ cơ thể theo người dùng \\
FastAPI / Uvicorn & 0.141.1 / 0.52.4 & Cung cấp REST API và WebSocket cho YOLO, đánh giá/gợi ý, thời tiết và nhận diện tư thế~\cite{FastAPIFeatures} \\
Ultralytics YOLO & 8.4.135 (backend) & Nạp \code{backend/models/best.pt}, phát hiện 21 lớp gốc và chuẩn hoá về tám nhóm dùng trong ứng dụng \\
PyTorch & 2.13.0 (backend); 2.11.0+cu130 (Colab) & Chạy suy luận trên CPU/CUDA và huấn luyện các phiên bản YOLO V1--V3 \\
Gemini / Google Gen AI & google-genai 2.25.0; Gemini 3.5 Flash-Lite & Đánh giá bộ đồ, sinh gợi ý có cấu trúc và trả lời hỏi đáp; không tạo kết quả giả khi dịch vụ lỗi \\
MediaPipe & 1.0.1 & Pose Landmarker xử lý khung hình camera và trả landmark/điểm neo cho thử đồ qua WebSocket~\cite{MediaPipePose} \\
Google Colab & GPU Tesla T4 & Chuẩn bị dữ liệu và huấn luyện mô hình YOLO trên môi trường trực tuyến \\
OpenCV / NumPy / Pillow & 5.0.0.93 / 2.5.2 / 12.3.0 & Giải mã, kiểm tra ảnh, trích màu và chuẩn bị đầu vào phía backend \\
Camera và xử lý ảnh Flutter & camera 0.12.1; image\_picker 1.2.3; image 4.9.2 & Chụp/chọn ảnh, chuyển đổi frame camera và ghép lớp trang phục \\
Kết nối và vị trí & http 1.6.0; WebSocket 3.0.3; geolocator 14.0.3 & Gọi REST API, truyền frame pose thời gian thực và lấy vị trí cho dữ liệu Open-Meteo \\
\bottomrule
\end{tabularx}
\end{table}

Phiên bản Ultralytics 8.4.135 ở Bảng~\ref{tab:congcu} là phiên bản trong
môi trường backend cục bộ. Notebook Colab không in phiên bản Ultralytics nên
báo cáo không đồng nhất hai môi trường này một cách suy đoán.

\section{Các công trình liên quan}

YOLO đặt bài toán phát hiện đối tượng thành một quy trình hồi quy thống nhất,
cho phép dự đoán hộp giới hạn và lớp trong một lần chạy mạng
~\cite{Redmon2016}. Đồ án sử dụng một checkpoint YOLO kích thước nano để ưu
tiên khả năng triển khai, sau đó fine-tune trên dữ liệu thời trang riêng.

DeepFashion cung cấp dữ liệu thời trang có chú thích phong phú cho các bài
toán nhận diện, truy hồi và phân tích thuộc tính quần áo
~\cite{DeepFashion2016}. Công trình cho thấy dữ liệu thời trang cần chú ý
đến biến thiên về tư thế, góc chụp, phông nền và sự khác biệt nhỏ giữa các
nhóm sản phẩm.

Han và cộng sự nghiên cứu khả năng tương thích thời trang bằng Bi-LSTM,
kết hợp đặc trưng ảnh với thông tin ngữ nghĩa để hỗ trợ gợi ý món đồ và toàn
bộ trang phục~\cite{Han2017}. Hướng tiếp cận này có khả năng học quan hệ từ
dữ liệu, trong khi đồ án hiện mô tả phạm vi gợi ý dựa trên loại, màu sắc và
phong cách cơ bản.

VITON tạo ảnh thử đồ bằng một quy trình học sâu từ thô đến tinh
~\cite{VITON2018}. Đồ án không triển khai mô hình sinh ảnh tương đương mà
chọn kỹ thuật ghép lớp 2D đơn giản hơn để phù hợp thời gian và tài nguyên.
Sự khác biệt này cần được nêu rõ khi đánh giá: chức năng thử đồ của sản phẩm
chỉ mang tính minh hoạ trực quan, không phải mô phỏng chân thực.

% ---------------------------------------------------------------------
%  CHƯƠNG 3 — PHÂN TÍCH VÀ THIẾT KẾ
% ---------------------------------------------------------------------
\chapter{Phân tích và thiết kế}
\label{ch:thietke}

\begin{muctieu}
Chương này phân tích đầu vào, đầu ra và yêu cầu của hệ thống; đồng thời mô
tả kiến trúc, dữ liệu và cách tổ chức các thành phần nhận diện, gợi ý phối
đồ và thử trang phục ảo.
\end{muctieu}

\section{Phân tích bài toán và yêu cầu}

Hệ thống hướng đến người dùng Android có nhu cầu lưu trữ quần áo cá nhân và
nhận hỗ trợ khi lựa chọn trang phục. Đầu vào chính là ảnh có một trang phục
nổi bật, ảnh người dùng phục vụ thử đồ và thông tin tài khoản. Đầu ra gồm
loại quần áo dự đoán, độ tin cậy, tối đa ba phương án phối đồ và ảnh xem
trước sau khi ghép lớp trang phục.

Các yêu cầu chức năng được xác định từ phạm vi ở Chương~\ref{ch:modau} gồm:
\begin{itemize}[leftmargin=1.6em, itemsep=3pt]
  \item đăng ký, đăng nhập và quản lý thông tin người dùng;
  \item nhận ảnh từ camera hoặc thư viện ảnh;
  \item nhận diện một số nhóm quần áo và trả độ tin cậy;
  \item thêm, sửa, xoá và xem các món đồ trong tủ đồ cá nhân;
  \item gợi ý tối đa ba phối đồ theo loại, màu sắc và phong cách;
  \item lưu trang phục, phối đồ hoặc kết quả yêu thích;
  \item ghép ảnh quần áo nền trong suốt lên ảnh người dùng và điều chỉnh vị
        trí, kích thước, góc xoay.
\end{itemize}

Yêu cầu phi chức năng gồm giao diện sử dụng được trên thiết bị Android, dữ
liệu mỗi người dùng được tách biệt, lỗi API được xử lý rõ ràng và thời gian
phản hồi hướng tới không quá 3 giây trong điều kiện kiểm thử ổn định.

Phạm vi tương tác giữa người dùng và các chức năng chính được tổng hợp ở
Hình~\ref{fig:usecase-tongquat}. Sơ đồ tập trung vào một tác nhân người dùng
và bảy use case đã xác định từ phạm vi đồ án; quan hệ với API và Firebase
được trình bày riêng trong sơ đồ kiến trúc ở Hình~\ref{fig:kientruc}.

\begin{figure}[H]
\centering
\resizebox{0.94\textwidth}{!}{%
\begin{tikzpicture}[
  uc/.style={draw=daumaroon, ellipse, align=center, fill=lightwarm,
             minimum width=3.5cm, minimum height=1.05cm, font=\small},
  assoc/.style={draw=gray700, line width=0.55pt}
]
  \draw[draw=daunavy, rounded corners=4pt, line width=0.8pt]
    (2.2,0.9) rectangle (12.0,-7.8);
  \node[anchor=west, font=\bfseries\color{daunavy}] at (2.45,0.55)
    {Hệ thống gợi ý phối đồ};

  % Tác nhân người dùng
  \draw[daunavy, line width=0.8pt] (0.7,-2.0) circle (0.28);
  \draw[daunavy, line width=0.8pt] (0.7,-2.28) -- (0.7,-3.25);
  \draw[daunavy, line width=0.8pt] (0.15,-2.62) -- (1.25,-2.62);
  \draw[daunavy, line width=0.8pt] (0.7,-3.25) -- (0.2,-4.0);
  \draw[daunavy, line width=0.8pt] (0.7,-3.25) -- (1.2,-4.0);
  \node[align=center, font=\small\bfseries, text=daunavy] at (0.7,-4.45)
    {Người dùng};

  \node[uc] (auth) at (4.5,-0.8) {Đăng ký, đăng nhập};
  \node[uc] (capture) at (9.3,-0.8) {Chọn hoặc chụp ảnh};
  \node[uc] (wardrobe) at (4.5,-2.65) {Quản lý tủ đồ};
  \node[uc] (recognize) at (9.3,-2.65) {Nhận diện trang phục};
  \node[uc] (recommend) at (4.5,-4.5) {Nhận gợi ý phối đồ};
  \node[uc] (favorite) at (9.3,-4.5) {Lưu và yêu thích};
  \node[uc] (tryon) at (6.9,-6.55) {Thử trang phục 2D};

  % Các đoạn trùng nhau được gom theo hai trục để tránh cắt qua ellipse.
  \draw[assoc] (1.25,-2.62) -- (2.65,-2.62);
  \draw[assoc] (2.65,-0.8) -- (2.65,-6.55);
  \draw[assoc] (2.65,-0.8) -- (auth.west);
  \draw[assoc] (2.65,-2.65) -- (wardrobe.west);
  \draw[assoc] (2.65,-4.5) -- (recommend.west);
  \draw[assoc] (2.65,-6.55) -- (tryon.west);

  \draw[assoc] (1.25,-2.62) -- (1.65,-2.62) -- (1.65,-7.45)
    -- (11.55,-7.45) -- (11.55,-0.8);
  \draw[assoc] (capture.east) -- (11.55,-0.8);
  \draw[assoc] (recognize.east) -- (11.55,-2.65);
  \draw[assoc] (favorite.east) -- (11.55,-4.5);
  \draw[assoc] (6.9,-7.45) -- (tryon.south);
\end{tikzpicture}}
\caption{Sơ đồ use case tổng quát của hệ thống gợi ý phối đồ.}
\label{fig:usecase-tongquat}
\end{figure}

Các use case trong sơ đồ được đặc tả ngắn gọn ở
Bảng~\ref{tab:usecase}. Đầu vào và kết quả chỉ mô tả những gì đã được xác
định trong phạm vi đồ án, không giả định tên màn hình hoặc endpoint cụ thể.

\begin{table}[H]
\centering
\caption{Danh sách use case chính của hệ thống.}
\label{tab:usecase}
\footnotesize
\begin{tabularx}{\textwidth}{>{\bfseries\color{daumaroon}}l >{\bfseries\color{daunavy}\raggedright\arraybackslash}p{3.0cm} >{\raggedright\arraybackslash}X >{\raggedright\arraybackslash}X}
\toprule
\rowcolor{gray100}\bfseries\color{daumaroon}Mã & \bfseries\color{daunavy}Use case & \bfseries\color{daumaroon}Đầu vào/thao tác & \bfseries\color{daunavy}Kết quả mong đợi \\
\midrule
UC01 & Xác thực & đăng ký hoặc đăng nhập & phiên người dùng và dữ liệu đúng tài khoản \\
UC02 & Tiếp nhận ảnh & chụp ảnh hoặc chọn từ thư viện & ảnh được kiểm tra và hiển thị xem trước \\
UC03 & Nhận diện trang phục & ảnh có một trang phục chính & nhãn dự đoán và độ tin cậy \\
UC04 & Quản lý tủ đồ & thêm, xem, sửa hoặc xoá món đồ & tủ đồ được cập nhật theo tài khoản \\
UC05 & Gợi ý phối đồ & món đồ đầu vào và dữ liệu tủ đồ & tối đa ba phương án không trùng \\
UC06 & Lưu và yêu thích & trang phục, phối đồ hoặc kết quả & nội dung được lưu để xem lại \\
UC07 & Thử đồ 2D & ảnh người dùng và ảnh trang phục & lớp ảnh có thể di chuyển, co giãn và xoay \\
\bottomrule
\end{tabularx}
\end{table}

Luồng xử lý chi tiết của bảy use case được minh hoạ lần lượt qua chức năng
xác thực ở Hình~\ref{fig:uc01-xac-thuc}, tiếp nhận ảnh ở
Hình~\ref{fig:uc02-tiep-nhan-anh}, nhận diện trang phục ở
Hình~\ref{fig:uc03-nhan-dien}, quản lý tủ đồ ở Hình~\ref{fig:uc04-tu-do},
gợi ý phối đồ ở Hình~\ref{fig:uc05-goi-y}, lưu nội dung và yêu thích ở
Hình~\ref{fig:uc06-luu-yeu-thich}, và thử trang phục 2D ở
Hình~\ref{fig:uc07-thu-do-2d}. Các sơ đồ cụ thể hoá những thao tác đã tổng
hợp trong Bảng~\ref{tab:usecase} mà không bổ sung chức năng ngoài phạm vi
đã xác định.

\begin{figure}[H]
\centering
\includegraphics[width=0.98\textwidth]{use-cases/uc01-xac-thuc.png}
\caption{Use case UC01: đăng ký và đăng nhập.}
\label{fig:uc01-xac-thuc}
\end{figure}

\begin{figure}[H]
\centering
\includegraphics[width=0.98\textwidth]{use-cases/uc02-tiep-nhan-anh.png}
\caption{Use case UC02: chọn hoặc chụp ảnh.}
\label{fig:uc02-tiep-nhan-anh}
\end{figure}

\begin{figure}[H]
\centering
\includegraphics[width=0.98\textwidth]{use-cases/uc03-nhan-dien.png}
\caption{Use case UC03: nhận diện trang phục.}
\label{fig:uc03-nhan-dien}
\end{figure}

\begin{figure}[H]
\centering
\includegraphics[width=0.98\textwidth]{use-cases/uc04-tu-do.png}
\caption{Use case UC04: quản lý tủ đồ.}
\label{fig:uc04-tu-do}
\end{figure}

\begin{figure}[H]
\centering
\includegraphics[width=0.98\textwidth]{use-cases/uc05-goi-y.png}
\caption{Use case UC05: gợi ý phối đồ.}
\label{fig:uc05-goi-y}
\end{figure}

\begin{figure}[H]
\centering
\includegraphics[width=0.98\textwidth]{use-cases/uc06-luu-yeu-thich.png}
\caption{Use case UC06: lưu nội dung và yêu thích.}
\label{fig:uc06-luu-yeu-thich}
\end{figure}

\begin{figure}[H]
\centering
\includegraphics[width=0.98\textwidth]{use-cases/uc07-thu-do-2d.png}
\caption{Use case UC07: thử trang phục 2D.}
\label{fig:uc07-thu-do-2d}
\end{figure}

Hệ thống hiện chỉ có một vai trò nghiệp vụ là người dùng đã xác thực. Firestore
cho phép đối chiếu \code{request.auth.uid} với định danh trên đường dẫn để giới
hạn người dùng vào dữ liệu của chính họ~\cite{FirestoreRules}. Quy tắc phiên bản
2 của dự án cho phép người dùng đọc, tạo, cập nhật và xoá đúng tài
liệu \code{users/\{userId\}} khi \code{request.auth.uid=userId}; quyền đọc/ghi
subcollection \code{wardrobe/\{itemId\}} cũng dùng cùng điều kiện. Người chưa
đăng nhập, truy cập UID khác hoặc truy cập đường dẫn không được khai báo đều
bị từ chối mặc định. Quy tắc chưa kiểm tra tên trường, kiểu dữ liệu hay giới
hạn nội dung, vì vậy người sở hữu vẫn có thể ghi mọi trường vào tài liệu của
mình; đây là điểm cần tăng cường trước khi triển khai thực tế.

Endpoint phân tích trang phục nhận multipart gồm ảnh JPG/JPEG/PNG tối đa
10 MB, hồ sơ, bối cảnh và ngưỡng confidence; phản hồi gồm các đối tượng
nhận diện và đánh giá Gemini nếu dịch vụ khả dụng. WebSocket thử đồ nhận cấu
hình camera trước, sau đó nhận JPEG nhị phân tối đa 3 MB và trả landmark/điểm
neo. Tiêu chí nghiệm thu là dữ liệu hợp lệ phải trả đúng schema; ảnh rỗng, sai
định dạng hoặc quá kích thước phải bị từ chối rõ ràng; tài khoản A không được
đọc/ghi dữ liệu của tài khoản B; ảnh và frame camera không được lưu tại backend.

\section{Thiết kế kiến trúc hệ thống}

Kiến trúc logic theo phạm vi đã mô tả được trình bày ở
Hình~\ref{fig:kientruc}. Sơ đồ chỉ thể hiện trách nhiệm của các khối; tên
dịch vụ, endpoint và vị trí triển khai phải được cập nhật theo mã nguồn.

\begin{figure}[H]
\centering
\resizebox{\textwidth}{!}{%
\begin{tikzpicture}[node distance=11mm and 14mm]
  \node[objbox] (user) {Người dùng\\Android};
  \node[clsbox, right=of user] (app) {Ứng dụng\\Flutter};
  \node[clsbox, right=of app] (api) {Dịch vụ\\API};
  \node[clsbox, right=of api] (ai) {YOLO, Gemini\\và MediaPipe};
  \node[objbox, right=of ai] (rec) {Kết quả phân tích,\\gợi ý và tư thế};
  \node[membox, below=of app] (overlay) {Ghép ảnh theo tư thế\\trên thiết bị};
  \node[membox, below=of api] (firebase) {Firebase Auth/Firestore\\tài khoản, hồ sơ, tủ đồ};
  \draw[flowarr] (user) -- (app);
  \draw[flowarr] (app) -- (api);
  \draw[flowarr] (api) -- (ai);
  \draw[flowarr] (ai) -- (rec);
  \draw[flowarr] (rec.south) -- ++(0,-8mm) -| (app.south);
  \draw[flowarr] (app) -- (overlay);
  \draw[flowarr] (app) -- (firebase);
\end{tikzpicture}}
\caption{Kiến trúc logic của hệ thống gợi ý phối đồ.}
\label{fig:kientruc}
\end{figure}

Ứng dụng Flutter chịu trách nhiệm tương tác, nhận ảnh, gọi API và hiển thị
kết quả. FastAPI là ranh giới với YOLO, Gemini, Open-Meteo và MediaPipe.
Flutter kết nối trực tiếp Firebase Authentication/Firestore; backend không
truy cập Firebase. Firestore đang lưu tài khoản và hồ sơ cơ thể, đồng thời
rules đã dành đường dẫn \code{users/\{uid\}/wardrobe/\{itemId\}} cho tủ đồ,
nhưng mã Flutter hiện chưa có luồng đọc/ghi subcollection này. Lớp trang phục
được Flutter ghép theo điểm neo tư thế và cho phép điều chỉnh thủ công.

\section{Thiết kế dữ liệu}

Từ mã nguồn ứng dụng, các nhóm dữ liệu thực tế được tổng hợp trong
Bảng~\ref{tab:dulieu}. Firebase dùng collection \code{users}; lịch sử hỏi
đáp và tuỳ chọn phân tích được lưu cục bộ theo định danh người dùng.

\begin{table}[H]
\centering
\caption{Các nhóm dữ liệu logic của hệ thống.}
\label{tab:dulieu}
\small
\begin{tabularx}{\textwidth}{>{\bfseries\color{daunavy}}l >{\raggedright\arraybackslash}X >{\raggedright\arraybackslash}X}
\toprule
\rowcolor{gray100}\bfseries\color{daunavy}Nhóm dữ liệu & \bfseries\color{daumaroon}Thông tin chính & \bfseries\color{daunavy}Mục đích \\
\midrule
Tài khoản & uid, họ tên, email, thời điểm tạo/cập nhật & xác thực và phân tách dữ liệu theo người dùng \\
Hồ sơ cơ thể & tên hồ sơ, chiều cao, cân nặng, giới tính tự khai, hồ sơ đang chọn & cá nhân hoá phân tích, gợi ý và hỏi đáp \\
Tuỳ chọn phân tích & chế độ thời tiết, mùa, nhiệt độ, phong cách, hoàn cảnh & khôi phục bối cảnh gần nhất theo tài khoản \\
Lịch sử hỏi đáp & vai trò, nội dung và thứ tự lượt chat & giữ hội thoại cục bộ và gửi tối đa sáu lượt gần nhất cho Gemini \\
Tủ đồ (đã có rules) & \code{users/\{uid\}/}\newline\code{wardrobe/\{itemId\}}; cấu trúc trường chưa xác minh & giới hạn quyền theo chủ tài khoản; mã Flutter hiện chưa sử dụng \\
Kết quả phân tích & lớp gốc, nhóm chuẩn hoá, màu, độ tin cậy, hộp và các điểm đánh giá & hiển thị kết quả YOLO/Gemini; không lưu ảnh ở backend \\
\bottomrule
\end{tabularx}
\end{table}

Ngoài dữ liệu nghiệp vụ, thành phần YOLO sử dụng hai nguồn ảnh. Bộ
\emph{Clothes Detection v1} xuất từ Roboflow gồm 11.926 ảnh, được chia thành
8.348 ảnh train, 2.385 ảnh validation và 1.193 ảnh test. File cấu hình công
bố giấy phép CC BY 4.0 và 21 lớp theo đúng thứ tự:
bag, belt, boat, coat, dress, glasses, hat, jacket, jumper, pants, pocket,
shirt, shoe, shoesClassic, shoesSports, shorts, skirt, sunglass, trousers,
tshirt và zipper.

Nguồn thứ hai là DeepFashion In-shop với 52.712 ảnh và hộp giới hạn
~\cite{DeepFashion2016}. Notebook V3 chuyển thành công 52.711 ảnh, bỏ qua một
ảnh lỗi, đồng thời ánh xạ 17 nhóm thư mục nguồn về taxonomy 21 lớp. Hai ánh
xạ gần đúng là \code{Rompers\_Jumpsuits} sang dress và \code{Suiting} sang
jacket. DeepFashion chỉ cho phép sử dụng phi thương mại trong nghiên cứu và
giáo dục theo README đi kèm dữ liệu.

Dữ liệu DeepFashion được chia theo \code{item\_id} với tỉ lệ 80/10/10 và
seed 42 để tránh rò rỉ nhiều góc chụp của cùng sản phẩm. Ảnh train của bộ dữ
liệu cũ được lặp hai lần trước khi trộn. Quy mô dữ liệu V3 sau khi chuẩn bị
được trình bày ở Bảng~\ref{tab:dulieu-yolo}.

\begin{table}[H]
\centering
\caption{Quy mô dữ liệu phát hiện trang phục ở phiên bản V3.}
\label{tab:dulieu-yolo}
\small
\begin{tabularx}{\textwidth}{>{\bfseries\color{daunavy}}l >{\raggedleft\arraybackslash}X >{\raggedleft\arraybackslash}X >{\raggedright\arraybackslash}p{6.0cm}}
\toprule
\rowcolor{gray100}\bfseries\color{daunavy}Tập & \bfseries\color{daumaroon}Ảnh & \bfseries\color{daumaroon}Đối tượng & \bfseries\color{daunavy}Ghi chú \\
\midrule
Train & 59.463 & 69.715 & gồm 16.696 lượt ảnh cũ sau khi lặp và 42.767 ảnh DeepFashion \\
Validation & 6.902 & 9.275 & gồm 2.385 ảnh cũ và 4.517 ảnh DeepFashion \\
Test & 6.620 & 7.821 & gồm 1.193 ảnh cũ và 5.427 ảnh DeepFashion \\
\midrule
Tổng & 72.985 & 86.811 & số ảnh train bao gồm lượt lặp có chủ đích \\
\bottomrule
\end{tabularx}
\end{table}

Nhật ký Ultralytics thực tế chấp nhận 59.315 ảnh train và 6.882 ảnh
validation; phần còn lại bị bỏ qua chủ yếu do nhãn cũ trộn hàng detection và
segmentation hoặc có toạ độ ngoài khoảng $[0,1]$. Đây là vấn đề chất lượng dữ
liệu cần được xử lý trước khi huấn luyện lại. Mã Flutter xác nhận hồ sơ tài
khoản và mảng \code{bodyProfiles} nằm trong \code{users/\{uid\}}; ảnh phân
tích, frame camera và kết quả YOLO không được ghi vào Firestore.

\section{Thiết kế mô hình hoặc cơ sở tri thức}

Thành phần nhận diện là bộ phát hiện đối tượng YOLO. Chuỗi xử lý gồm kiểm tra
ảnh, đưa ảnh về kích thước 640 pixel, suy luận hộp giới hạn và điểm lớp, lọc
kết quả theo ngưỡng của dịch vụ rồi trả nhãn, độ tin cậy và vị trí trang phục.
Các notebook và checkpoint cho thấy ba vòng lặp huấn luyện trong
Bảng~\ref{tab:cauhinh-yolo}; đây là các phiên bản fine-tune của cùng họ mô
hình chứ không phải ba kiến trúc độc lập.

\begin{table}[H]
\centering
\caption{Cấu hình các phiên bản huấn luyện YOLO.}
\label{tab:cauhinh-yolo}
\footnotesize
\begin{tabularx}{\textwidth}{>{\bfseries\color{daunavy}}l >{\raggedright\arraybackslash}p{2.4cm} >{\raggedright\arraybackslash}X >{\raggedright\arraybackslash}p{4.2cm}}
\toprule
\rowcolor{gray100}\bfseries\color{daunavy}Phiên bản & \bfseries\color{daumaroon}Khởi tạo & \bfseries\color{daunavy}Dữ liệu & \bfseries\color{daumaroon}Cấu hình chính \\
\midrule
V1 & \code{yolo26n.pt} & Roboflow 11.926 ảnh & tối đa 60 epoch, batch 16, ảnh 640, optimizer tự động \\
V2 & \code{best.pt} & dữ liệu cũ và 6.500 ảnh DeepFashion được chọn & tối đa 40 epoch, patience 10, AdamW, $lr_0=10^{-4}$, cosine LR \\
V3 & checkpoint V1/V3 & dữ liệu ở Bảng~\ref{tab:dulieu-yolo} & 30 epoch đã ghi nhận và 5 epoch tiếp tục, AdamW, batch 16, ảnh 640, seed 42 \\
\bottomrule
\end{tabularx}
\end{table}

Ở V3, lần chạy chính dùng $lr_0=10^{-4}$, weight decay 0,0005, cosine LR,
mosaic 0,5 và AMP. Sau khi phiên chạy dừng ở epoch 30, notebook nạp
\code{last.pt} và chạy tiếp 5 epoch với $lr_0=10^{-5}$, tắt AMP và lưu
checkpoint \code{best.pt}. Các tệp kết quả ghi nhận box loss, classification
loss và L1 loss cho cả train và validation.

Thành phần phân tích gửi metadata nhận diện, hồ sơ và bối cảnh đến Gemini để
nhận các điểm màu sắc, mùa/nhiệt độ, tỉ lệ cơ thể, phong cách và hoàn cảnh.
Điểm tổng dùng trọng số tương ứng 0,30; 0,30; 0,20; 0,10 và 0,10; tiêu chí
thiếu dữ liệu được bỏ khỏi mẫu số thay vì tính bằng 0. Luồng Gợi ý độc lập
với ảnh và YOLO, trả một nhận định, một lý do và từ một đến ba điều chỉnh.
Khi Gemini lỗi, hệ thống giữ kết quả YOLO và không tạo đánh giá giả.

Thành phần thử đồ nhận JPEG camera, chạy MediaPipe Pose ở chế độ
\code{VIDEO}, trả landmark và điểm neo qua WebSocket. Flutter ánh xạ toạ độ
camera, làm mượt tư thế rồi đặt lớp ảnh theo vai, hông hoặc bàn chân; người
dùng có thể tinh chỉnh vị trí, tỉ lệ và góc xoay. Phương pháp này không suy
diễn chính xác hình dạng ba chiều của cơ thể hoặc đặc tính vật lý của vải.

API mặc định nạp \code{backend/models/best.pt}; endpoint phân tích dùng
confidence 0,25 và IoU 0,45. Trên giao diện, người dùng có thể điều chỉnh
confidence trong khoảng 0,10--0,90. Các thông tin này được lấy trực tiếp từ
mã backend và mã Flutter.

% ---------------------------------------------------------------------
%  CHƯƠNG 4 — CÀI ĐẶT VÀ THỬ NGHIỆM
% ---------------------------------------------------------------------
\chapter{Cài đặt và thử nghiệm}
\label{ch:caidat}

\begin{muctieu}
Chương này mô tả môi trường triển khai, trách nhiệm của các thành phần, giao
diện đã cài đặt và kết quả kiểm thử hệ thống.
\end{muctieu}

\section{Môi trường triển khai}

Ứng dụng \code{app\_nhan\_dien} dùng Flutter 3.47.1 stable và Dart 3.13.1.
Mã nguồn có \code{FirebaseOptions} cho Android, iOS, Web và Windows; phạm vi
cài đặt, chạy thử và kiểm thử của báo cáo tập trung vào Android. Firebase
Authentication hỗ trợ xác thực email/mật khẩu~\cite{FirebaseAuthFlutter}.
Cloud Firestore tổ chức dữ liệu theo tài liệu và collection~\cite{FirestoreDataModel};
trong mã dự án, dịch vụ này lưu tài liệu \code{users/\{uid\}} và trường hồ sơ cơ thể. Các tuỳ chọn phân tích và lịch
sử hỏi đáp được lưu cục bộ bằng \code{shared\_preferences} 2.5.5. Thành phần
trí tuệ nhân tạo được truy cập qua FastAPI 0.141.1 chạy bằng Uvicorn 0.52.4.

Thông tin trên Firebase Console được đối chiếu với cấu hình và lời gọi dịch
vụ trong mã nguồn, kết quả trình bày ở Bảng~\ref{tab:cauhinh-firebase}. Cách
đối chiếu này tách biệt dịch vụ đã đăng ký trên Console với dịch vụ thật sự
được ứng dụng gọi trong phiên bản hiện tại.

\begin{table}[H]
\centering
\caption{Cấu hình dự án Firebase tại thời điểm đối chiếu ngày 11/11/2026.}
\label{tab:cauhinh-firebase}
\footnotesize
\begin{tabularx}{\textwidth}{>{\bfseries\color{daunavy}\raggedright\arraybackslash}p{3.4cm} >{\raggedright\arraybackslash}X}
\toprule
\rowcolor{gray100}\bfseries\color{daunavy}Hạng mục & \bfseries\color{daumaroon}Cấu hình và trạng thái sử dụng \\
\midrule
Dự án và gói dịch vụ & Tên hiển thị \code{ai-fashion-project}, Project ID \code{ai-fashion-project-6fa98}; gói Spark (no-cost). \\
Ứng dụng đã đăng ký & Android \code{com.example.app\_nhan\_dien}, iOS \code{com.example.appNhanDien} và hai cấu hình Web. Các định danh này khớp với \code{firebase\_options.dart} và \code{google-services.json}. \\
Authentication & Email/Password đã bật. Mã Flutter thực hiện đăng ký, đăng nhập, đăng xuất, đặt lại và đổi mật khẩu bằng Firebase Auth. Tên miền được uỷ quyền gồm \code{localhost}, \code{ai-fashion-project-6fa98.firebaseapp.com} và \code{ai-fashion-project-6fa98.web.app}. \\
Cloud Firestore & Cơ sở dữ liệu Native đặt tại \code{nam5}. Ứng dụng đọc/ghi hồ sơ tại \code{users/\{uid\}}; Rules chỉ cho người đã xác thực thao tác tài liệu của chính mình và subcollection \code{wardrobe}. \\
Google Analytics & Firebase Console đã liên kết Property ID \code{553061587}; cấu hình Web có Measurement ID. Tuy nhiên, mã Flutter không khai báo \code{firebase\_analytics} và chưa có lời gọi ghi sự kiện, nên báo cáo không xem Analytics là chức năng đã triển khai. \\
Firebase AI Logic & Global Context đã được cấu hình ở vùng \code{global} trên Console. Phiên bản hiện tại chưa dùng Firebase AI SDK; Flutter gửi HTTP đến FastAPI, sau đó backend gọi Gemini bằng thư viện \code{google-genai}. \\
\bottomrule
\end{tabularx}
\end{table}

Môi trường huấn luyện được ghi nhận từ notebook Google Colab: một GPU Tesla
T4, PyTorch 2.11.0+cu130 và CUDA khả dụng. Dữ liệu được sao chép từ Google
Drive sang \code{/content} trước khi huấn luyện để giảm độ trễ đọc hàng nghìn
tệp nhỏ. Môi trường triển khai cục bộ dùng Python 3.13.7, Ultralytics 8.4.135,
PyTorch 2.13.0 và MediaPipe 1.0.1; model mặc định nằm tại
\code{backend/models/best.pt}.

Các biến cấu hình quan trọng gồm \code{MODEL\_PATH},
\code{POSE\_MODEL\_PATH}, \code{GEMINI\_API\_KEY},
\code{GEMINI\_MODEL} và \code{GEMINI\_TIMEOUT\_SECONDS}. Khóa API không được
ghi vào mã nguồn. Báo cáo chưa thực hiện benchmark trên một cấu hình cố định
gồm CPU, RAM, Android SDK và thiết bị thật, vì vậy không công bố số liệu thời
gian phản hồi theo phần cứng.

\section{Các thành phần chính}

Ứng dụng di động chịu trách nhiệm xác thực, quản lý hồ sơ cơ thể, nhận ảnh,
gửi yêu cầu phân tích, lấy gợi ý, hỏi đáp và hiển thị thử đồ theo camera.
Giao diện chính có bốn tab: Phân tích, Gợi ý, Hỏi đáp và Cá nhân. Các thao
tác mạng có trạng thái đang xử lý và lỗi; lịch sử hỏi đáp được giữ theo tài
khoản ở bộ nhớ cục bộ.

Dịch vụ API tiếp nhận JPG/PNG tối đa 10 MB, kiểm tra MIME và nội dung ảnh,
sau đó gọi checkpoint YOLO 21 lớp. Kết quả giữ nhiều hộp hợp lệ, gồm lớp gốc,
nhóm chuẩn hoá, màu chính/phụ, độ tin cậy và toạ độ. Endpoint phân tích dùng
confidence mặc định 0,25 và IoU 0,45, rồi gửi metadata cùng hồ sơ/bối cảnh
đến Gemini. Các endpoint gợi ý và hỏi đáp là luồng không ảnh, không dùng kết
quả YOLO hay dữ liệu tủ đồ.

Firebase Authentication quản lý phiên đăng nhập. Firestore lưu hồ sơ tài
khoản và danh sách hồ sơ cơ thể trong tài liệu \code{users/\{uid\}}. Security
Rules cho phép chủ tài khoản đọc/ghi subcollection
\code{wardrobe/\{itemId\}}, nhưng mã Flutter hiện chưa triển khai màn hình và
service tủ đồ; không có kho ảnh hoặc danh sách yêu thích. Dữ liệu ảnh gửi đến
backend chỉ tồn tại trong request và không được ghi ra ổ đĩa.

MediaPipe Pose Landmarker cung cấp các điểm mốc cơ thể từ ảnh hoặc video
~\cite{MediaPipePose}. Trong hệ thống, chức năng thử đồ gửi frame JPEG qua
WebSocket đến Pose Landmarker. Backend trả landmark và điểm neo; Flutter làm
mượt tư thế rồi ghép lớp trang
phục theo vùng thân trên, thân dưới, váy hoặc bàn chân. Người dùng có thể dịch
chuyển, co giãn, xoay và đặt lại lớp ảnh; frame không được lưu ở backend.

Mã huấn luyện xác nhận quy trình nạp \code{yolo26n.pt}, chuyển annotation
DeepFashion sang YOLO, kiểm tra thứ tự 21 lớp, fine-tune và lưu
\code{best.pt}. Mã triển khai nạp checkpoint này một lần khi backend khởi
động, tự chọn \code{cuda:0} nếu CUDA khả dụng và dùng CPU trong trường hợp
còn lại. Nếu Gemini thiếu khóa, hết hạn mức hoặc lỗi mạng, API vẫn trả kết
quả YOLO với \code{source=unavailable} thay vì tạo đánh giá thay thế giả.

\section{Giao diện và demo}

Luồng giao diện thực tế gồm đăng ký/đăng nhập, phân tích ảnh từ camera hoặc
thư viện, gợi ý theo bối cảnh, hỏi đáp tiếp nối, quản lý hồ sơ cá nhân và thử
trang phục theo camera. Kết quả phân tích hiển thị nhãn, màu, độ tin cậy,
điểm đánh giá và đề xuất điều chỉnh khi Gemini khả dụng.

Màn hình xác thực trong Hình~\ref{fig:ui-auth} cung cấp hai luồng tạo tài
khoản và đăng nhập bằng email/mật khẩu. Sau khi đăng nhập, người dùng truy
cập bốn nhóm chức năng chính bằng thanh điều hướng dưới màn hình.

\begin{figure}[H]
\centering
\begin{minipage}[t]{0.46\textwidth}
  \centering
  \includegraphics[width=\linewidth,height=0.57\textheight,keepaspectratio]{screenshots/ui-login.jpg}

  \small (a) Đăng nhập tài khoản
\end{minipage}\hfill
\begin{minipage}[t]{0.46\textwidth}
  \centering
  \includegraphics[width=\linewidth,height=0.57\textheight,keepaspectratio]{screenshots/ui-register.jpg}

  \small (b) Đăng ký tài khoản
\end{minipage}
\caption{Giao diện xác thực người dùng bằng email và mật khẩu.}
\label{fig:ui-auth}
\end{figure}

Khi quên mật khẩu, người dùng nhập email đã đăng ký và yêu cầu Firebase gửi
liên kết đặt lại. Hình~\ref{fig:ui-password-reset} thể hiện hộp nhập email và
thông báo xác nhận sau khi yêu cầu được gửi thành công.

\begin{figure}[H]
\centering
\begin{minipage}[t]{0.46\textwidth}
  \centering
  \includegraphics[width=\linewidth,height=0.57\textheight,keepaspectratio]{screenshots/ui-password-reset-dialog.jpg}

  \small (a) Nhập email đặt lại mật khẩu
\end{minipage}\hfill
\begin{minipage}[t]{0.46\textwidth}
  \centering
  \includegraphics[width=\linewidth,height=0.57\textheight,keepaspectratio]{screenshots/ui-password-reset-sent.jpg}

  \small (b) Xác nhận đã gửi email
\end{minipage}
\caption{Luồng yêu cầu đặt lại mật khẩu qua email.}
\label{fig:ui-password-reset}
\end{figure}

Đối với tài khoản đang đăng nhập, chức năng đổi mật khẩu yêu cầu nhập lại mật
khẩu hiện tại trước khi lưu mật khẩu mới. Hình~\ref{fig:ui-password-change}
minh hoạ biểu mẫu đổi mật khẩu và thông báo hoàn tất trên màn hình tài khoản.

\begin{figure}[H]
\centering
\begin{minipage}[t]{0.46\textwidth}
  \centering
  \includegraphics[width=\linewidth,height=0.57\textheight,keepaspectratio]{screenshots/ui-password-change-dialog.jpg}

  \small (a) Nhập mật khẩu hiện tại và mật khẩu mới
\end{minipage}\hfill
\begin{minipage}[t]{0.46\textwidth}
  \centering
  \includegraphics[width=\linewidth,height=0.57\textheight,keepaspectratio]{screenshots/ui-password-change-success.jpg}

  \small (b) Thông báo đổi mật khẩu thành công
\end{minipage}
\caption{Luồng đổi mật khẩu của tài khoản đang đăng nhập.}
\label{fig:ui-password-change}
\end{figure}

Luồng phân tích bắt đầu từ màn hình chọn ảnh hoặc chụp ảnh. Sau khi backend
xử lý, ứng dụng vẽ hộp giới hạn cùng lớp trang phục, màu và độ tin cậy như
Hình~\ref{fig:ui-detection}. Đây là bằng chứng giao diện cho thấy kết quả YOLO
đã được trả về và hiển thị trên ảnh đầu vào.

\begin{figure}[H]
\centering
\begin{minipage}[t]{0.46\textwidth}
  \centering
  \includegraphics[width=\linewidth,height=0.57\textheight,keepaspectratio]{screenshots/ui-analysis-input.jpg}

  \small (a) Chọn nguồn ảnh và hồ sơ
\end{minipage}\hfill
\begin{minipage}[t]{0.46\textwidth}
  \centering
  \includegraphics[width=\linewidth,height=0.57\textheight,keepaspectratio]{screenshots/ui-analysis-detection.jpg}

  \small (b) Hiển thị kết quả nhận diện
\end{minipage}
\caption{Luồng tiếp nhận ảnh và nhận diện trang phục bằng YOLO.}
\label{fig:ui-detection}
\end{figure}

Phần đánh giá sau nhận diện được minh hoạ ở
Hình~\ref{fig:ui-analysis-result}. Ứng dụng trình bày điểm tổng hợp, các tiêu
chí màu sắc, thời tiết, phong cách và hoàn cảnh; đồng thời nêu điểm mạnh,
điểm cần điều chỉnh và khuyến nghị phối bổ sung. Trường hợp trong ảnh cũng
cho thấy hệ thống phát hiện sự không khớp giữa nhãn chân váy và hồ sơ giới
tính Nam thay vì bỏ qua bối cảnh người dùng.

\begin{figure}[H]
\centering
\begin{minipage}[t]{0.46\textwidth}
  \centering
  \includegraphics[width=\linewidth,height=0.57\textheight,keepaspectratio]{screenshots/ui-analysis-score.jpg}

  \small (a) Điểm và tiêu chí đánh giá
\end{minipage}\hfill
\begin{minipage}[t]{0.46\textwidth}
  \centering
  \includegraphics[width=\linewidth,height=0.57\textheight,keepaspectratio]{screenshots/ui-analysis-recommendations.jpg}

  \small (b) Điểm cần điều chỉnh và khuyến nghị
\end{minipage}
\caption{Kết quả đánh giá và khuyến nghị sau khi phân tích trang phục.}
\label{fig:ui-analysis-result}
\end{figure}

Với luồng gợi ý không cần ảnh, người dùng chọn hồ sơ, nguồn thời tiết, phong
cách và hoàn cảnh. Hình~\ref{fig:ui-suggestion} cho thấy dữ liệu Open-Meteo
được hiển thị trước khi gửi yêu cầu và kết quả Gemini trả về dưới dạng mô tả
bộ đồ cùng các cách điều chỉnh phù hợp.

\begin{figure}[H]
\centering
\begin{minipage}[t]{0.46\textwidth}
  \centering
  \includegraphics[width=\linewidth,height=0.57\textheight,keepaspectratio]{screenshots/ui-suggestion-context.jpg}

  \small (a) Chọn hồ sơ và bối cảnh thời tiết
\end{minipage}\hfill
\begin{minipage}[t]{0.46\textwidth}
  \centering
  \includegraphics[width=\linewidth,height=0.57\textheight,keepaspectratio]{screenshots/ui-suggestion-result.jpg}

  \small (b) Bộ đồ và cách điều chỉnh được đề xuất
\end{minipage}
\caption{Gợi ý phối đồ dựa trên hồ sơ, thời tiết và hoàn cảnh.}
\label{fig:ui-suggestion}
\end{figure}

Hình~\ref{fig:ui-chat} minh hoạ khả năng hỏi đáp tiếp nối: câu trả lời đầu
tiên sử dụng bối cảnh mưa và nhiệt độ hiện tại, còn lượt hỏi sau tiếp tục dùng
thông tin hồ sơ và lịch sử hội thoại để tư vấn loại giày phù hợp.

\begin{figure}[H]
\centering
\begin{minipage}[t]{0.46\textwidth}
  \centering
  \includegraphics[width=\linewidth,height=0.57\textheight,keepaspectratio]{screenshots/ui-chat-initial.jpg}

  \small (a) Câu hỏi và phản hồi đầu tiên
\end{minipage}\hfill
\begin{minipage}[t]{0.46\textwidth}
  \centering
  \includegraphics[width=\linewidth,height=0.57\textheight,keepaspectratio]{screenshots/ui-chat-followup.jpg}

  \small (b) Hỏi tiếp trong cùng ngữ cảnh
\end{minipage}
\caption{Giao diện hỏi đáp phối đồ nhiều lượt với Gemini.}
\label{fig:ui-chat}
\end{figure}

Người dùng có thể tạo nhiều hồ sơ cơ thể và thay đổi hồ sơ đang hoạt động.
Các biểu mẫu trong Hình~\ref{fig:ui-profile} thu thập tên hồ sơ, chiều cao,
cân nặng và giới tính; các dữ liệu này được dùng làm bối cảnh cho phân tích và
gợi ý phối đồ.

\begin{figure}[H]
\centering
\begin{minipage}[t]{0.46\textwidth}
  \centering
  \includegraphics[width=\linewidth,height=0.57\textheight,keepaspectratio]{screenshots/ui-profile-edit.jpg}

  \small (a) Chỉnh sửa hồ sơ cơ thể
\end{minipage}\hfill
\begin{minipage}[t]{0.46\textwidth}
  \centering
  \includegraphics[width=\linewidth,height=0.57\textheight,keepaspectratio]{screenshots/ui-profile-add.jpg}

  \small (b) Tạo hồ sơ cơ thể mới
\end{minipage}
\caption{Quản lý các hồ sơ cơ thể của người dùng.}
\label{fig:ui-profile}
\end{figure}

Chức năng thử đồ cho phép người dùng căn ảnh trang phục vào khung, phóng thu,
xoay, lật và xem trước trước khi tạo lớp ghép. Hình~\ref{fig:ui-tryon-prepare}
minh hoạ bước tạo filter và kết quả ghép áo 2D theo các điểm neo tư thế. Dòng
trạng thái kết nối và FPS đồng thời cho biết frame camera đang được xử lý theo
thời gian thực.

\begin{figure}[H]
\centering
\begin{minipage}[t]{0.46\textwidth}
  \centering
  \includegraphics[width=\linewidth,height=0.57\textheight,keepaspectratio]{screenshots/ui-tryon-calibrate.png}

  \small (a) Căn ảnh trang phục và tạo filter
\end{minipage}\hfill
\begin{minipage}[t]{0.46\textwidth}
  \centering
  \includegraphics[width=\linewidth,height=0.57\textheight,keepaspectratio]{screenshots/ui-tryon-2d.png}

  \small (b) Ghép lớp áo 2D theo tư thế
\end{minipage}
\caption{Chuẩn bị ảnh trang phục và thử đồ trực tiếp ở chế độ 2D.}
\label{fig:ui-tryon-prepare}
\end{figure}

Ngoài lớp ghép 2D, ứng dụng còn cung cấp chế độ 3D bám thân. Ảnh thực nghiệm
trong Hình~\ref{fig:ui-tryon-3d} cho thấy mesh áo được biến dạng theo phần cơ
thể nhìn thấy, đồng thời người dùng vẫn có thể chuyển về 2D, xoá lớp ảnh,
chọn nhóm trang phục khác hoặc thêm ảnh đồ mới. Camera chỉ được xử lý tạm
thời và không lưu frame tại backend.

\begin{figure}[H]
\centering
\includegraphics[width=0.52\textwidth,height=0.68\textheight,keepaspectratio]{screenshots/ui-tryon-3d.png}
\caption{Kết quả thử đồ trực tiếp ở chế độ 3D bám thân.}
\label{fig:ui-tryon-3d}
\end{figure}

\clearpage

\section{Kiểm thử}

Các bộ kiểm thử tự động được chạy lại ngày 11/11/2026 và tổng hợp ở
Bảng~\ref{tab:kiemthu}. Backend dùng mock cho YOLO, Gemini, MediaPipe và API
ngoài nên không phát sinh chi phí dịch vụ; kết quả Flutter phản ánh đúng trạng
thái mã tại thời điểm biên soạn báo cáo.

\begin{table}[H]
\centering
\caption{Các nhóm ca kiểm thử của hệ thống.}
\label{tab:kiemthu}
\footnotesize
\begin{tabularx}{\textwidth}{>{\bfseries\color{daunavy}}l >{\raggedright\arraybackslash}X >{\raggedright\arraybackslash}X >{\color{daumaroon}}l}
\toprule
\rowcolor{gray100}\bfseries\color{daunavy}Nhóm & \bfseries\color{daumaroon}Đầu vào/thao tác & \bfseries\color{daunavy}Kết quả mong đợi & \bfseries\color{daumaroon}Thực tế \\
\midrule
Backend & validation ảnh, YOLO, Gemini, thời tiết, pose và WebSocket & route trả dữ liệu/lỗi đúng schema, không lưu frame & 104/104 đạt \\
Flutter & model, service, widget và điều hướng ứng dụng & toàn bộ test tự động đạt & 139 đạt, 6 lỗi \\
Lỗi giao diện & hồ sơ cá nhân, gợi ý không ảnh, đổi giới tính, thời tiết và chuyển tab & nội dung và trạng thái khớp kỳ vọng test & chưa đạt \\
Thử đồ & ánh xạ camera, làm mượt pose, filter, mesh và thao tác lớp ảnh & tính toán ổn định trên test tự động & đạt trong bộ test \\
\bottomrule
\end{tabularx}
\end{table}

Sáu lỗi Flutter tập trung ở các test widget đã lệch với giao diện hiện tại:
hiển thị tên tại màn Cá nhân; trạng thái giới tính khi yêu cầu gợi ý đang chờ;
ba trường hợp nội dung Gợi ý/Open-Meteo; và luồng mở Gợi ý từ thanh tab. Đây
là lỗi cần đồng bộ giữa giao diện và test trước khi tuyên bố toàn bộ ứng dụng
đạt kiểm thử hồi quy.


% ---------------------------------------------------------------------
%  CHƯƠNG 5 — KẾT QUẢ VÀ ĐÁNH GIÁ
% ---------------------------------------------------------------------
\chapter{Kết quả và đánh giá}
\label{ch:ketqua}

\begin{muctieu}
Chương này đối chiếu các mục tiêu của đồ án, xác định độ đo đánh giá và nêu
những giới hạn có thể khẳng định từ phạm vi đã chọn.
\end{muctieu}

\section{Kết quả đạt được}

Thư mục \code{AI\_Fashion\_Project} cung cấp notebook, checkpoint, cấu hình
và nhật ký cho các lần chạy V1, V2 và V3. Nhờ đó có thể xác nhận thành phần
nhận diện đã được huấn luyện theo bài toán phát hiện đối tượng 21 lớp và đã
lưu các checkpoint \code{best.pt}. V2 có kết quả trên tập test riêng; V3 có
kết quả validation và các biểu đồ do Ultralytics sinh ra.

Mục tiêu ``độ chính xác 80\%'' ban đầu không tương ứng trực tiếp với chỉ số
của bộ phát hiện. Chỉ số tốt nhất ghi nhận cho V3 full là mAP@0,5:0,95 bằng
36,69\% trên validation, nên chưa có bằng chứng để tuyên bố đạt mục tiêu
80\%. Bộ kiểm thử backend đạt 104/104; bộ Flutter có 139 test đạt và 6 test
widget lỗi. Phạm vi thực nghiệm hiện tại chưa bao gồm phép đo thời gian API
trên thiết bị Android thật.

\section{Độ đo đánh giá}

Đối với bộ phát hiện đối tượng, precision và recall được tính theo
Công thức~\eqref{eq:precision-recall}, trong đó TP, FP và FN lần lượt là dự
đoán đúng, dự đoán sai và đối tượng bị bỏ sót:
\begin{equation}
\mathrm{Precision}=\frac{TP}{TP+FP},\qquad
\mathrm{Recall}=\frac{TP}{TP+FN}.
\label{eq:precision-recall}
\end{equation}

Average Precision là diện tích dưới đường cong precision--recall của một
lớp. mAP là trung bình AP trên các lớp; báo cáo dùng mAP@0,5 và
mAP@0,5:0,95 theo đúng tệp \code{results.csv} của Ultralytics.

Ngoài chất lượng nhận diện, thời gian từ lúc gửi yêu cầu đến khi nhận kết quả
là độ đo hiệu năng phù hợp cho hệ thống. Chức năng gợi ý có thể được đánh giá
qua tỉ lệ phương án hợp lệ, số kết quả trùng và đánh giá thủ công của người
dùng; chức năng thử đồ 2D được xem xét theo khả năng thao tác và độ ổn định
của lớp ảnh. Các độ đo này chưa được thu thập thành một bộ số liệu định lượng
trong phạm vi thực nghiệm của báo cáo.

Bảng~\ref{tab:ketqua} tổng hợp trực tiếp các tệp \code{results.csv} và đầu
ra đánh giá trong notebook. Các dòng không cùng tập dữ liệu, vì vậy chỉ phản
ánh tiến trình thực nghiệm và không được dùng như một so sánh tuyệt đối giữa
các kiến trúc.

\begin{table}[H]
\centering
\caption{Kết quả được ghi nhận ở các phiên bản YOLO.}
\label{tab:ketqua}
\footnotesize
\begin{tabularx}{\textwidth}{>{\bfseries\color{daunavy}}l >{\raggedright\arraybackslash}X c c c c}
\toprule
\rowcolor{gray100}\bfseries\color{daunavy}Lần chạy & \bfseries\color{daumaroon}Tập đánh giá & \bfseries\color{daunavy}Precision & \bfseries\color{daunavy}Recall & \bfseries\color{daumaroon}mAP@0,5 & \bfseries\color{daumaroon}mAP@0,5:0,95 \\
\midrule
V1, epoch 1 & Val. & 75,55\% & 26,56\% & 23,41\% & 15,98\% \\
V2, checkpoint tốt nhất & test V2 & 54,46\% & 39,95\% & 42,33\% & 30,01\% \\
V3 full, epoch 19 & Val. & 60,31\% & 51,03\% & 50,11\% & 36,69\% \\
V3, 5 epoch tiếp tục & Val. & 55,18\% & 49,19\% & 49,53\% & 36,10\% \\
\bottomrule
\end{tabularx}
\end{table}

Đường cong loss và các độ đo của năm epoch tiếp tục được trình bày ở
Hình~\ref{fig:v3-results}. Ma trận nhầm lẫn chuẩn hoá tương ứng ở
Hình~\ref{fig:v3-confusion}, còn một lô dự đoán validation được minh hoạ ở
Hình~\ref{fig:v3-predictions}.

\begin{figure}[H]
\centering
\includegraphics[width=\textwidth]{ai-training/v3-results.png}
\caption{Loss và độ đo của năm epoch fine-tune cuối ở V3.}
\label{fig:v3-results}
\end{figure}

\begin{figure}[H]
\centering
\includegraphics[width=0.92\textwidth]{ai-training/v3-confusion-matrix-normalized.png}
\caption{Ma trận nhầm lẫn chuẩn hoá của V3 trên tập validation.}
\label{fig:v3-confusion}
\end{figure}

\begin{figure}[H]
\centering
\includegraphics[width=\textwidth]{ai-training/v3-validation-predictions.jpg}
\caption{Một lô dự đoán của V3 trên tập validation.}
\label{fig:v3-predictions}
\end{figure}

Kết quả V3 trong báo cáo được lấy từ tập validation vì notebook chưa lưu đầu
ra đánh giá trên tập test kết hợp. Do đó, Bảng~\ref{tab:ketqua} chỉ phản ánh
các kết quả đã được ghi nhận và không suy rộng sang hiệu năng trên tập test
hoặc thời gian phản hồi của thiết bị triển khai.

\section{Bàn luận}

V2 cải thiện rõ mAP so với kết quả V1 được lưu, nhưng V1 mới chỉ có một epoch
trong \code{results.csv} nên phép so sánh này không đại diện cho tiềm năng
của cấu hình 60 epoch. Với V3 full, mAP@0,5:0,95 tốt nhất đạt 36,69\% ở epoch
19. Năm epoch tiếp tục làm train loss giảm, nhưng mAP không vượt kết quả tốt
nhất trước đó; điều này cho thấy việc kéo dài huấn luyện từ \code{last.pt}
không tự động cải thiện khả năng tổng quát hoá.

Ma trận ở Hình~\ref{fig:v3-confusion} cho thấy một số lớp có đường chéo cao,
chẳng hạn hat khoảng 0,86, shoe và sunglass khoảng 0,88. Ngược lại, coat,
belt, trousers, shoesClassic và shoesSports có tỉ lệ đúng thấp; glasses và
zipper có quá ít mẫu validation để đưa ra kết luận ổn định. Hình
~\ref{fig:v3-predictions} cũng cho thấy mô hình có thể phát hiện đúng áo,
giày và quần trên ảnh sản phẩm, nhưng vẫn sinh hộp sai trên ảnh quảng cáo,
ảnh ghép và trường hợp có nhiều món đồ.

Với gợi ý phối đồ, kết quả do Gemini sinh từ hồ sơ và bối cảnh, không phải
thuật toán lấy tối đa ba bộ đồ trong tủ đồ. Báo cáo mới xác nhận luồng xử lý
và kết quả minh hoạ, chưa có khảo sát người dùng để định lượng tính phù hợp
và tính ổn định của phản hồi. Với thử đồ, các test tự động đã kiểm tra ánh xạ
và biến đổi; sai lệch vị trí trên camera thật khi tư thế, ánh sáng và khoảng
cách thay đổi chưa được lượng hoá.

\section{Hạn chế}

Phạm vi ở Chương~\ref{ch:modau} cho thấy hệ thống chỉ nhận diện một trang
phục chính và chưa xác định thương hiệu, chất liệu hoặc kích cỡ chính xác.
Chức năng phân tích dựa trên metadata YOLO và bối cảnh; chức năng Gợi ý không
dùng ảnh, không dùng tủ đồ và phụ thuộc vào Gemini. Hệ thống chưa học sở thích
cá nhân từ phản hồi dài hạn và chưa có bộ dữ liệu chuẩn để chấm chất lượng tư
vấn.

Thử đồ bằng lớp ảnh 2D không tái tạo độ rủ, nếp gấp, che khuất, ánh sáng hay
biến dạng theo cơ thể. Người dùng vẫn phải điều chỉnh vị trí, kích thước và
góc xoay. Ứng dụng chủ yếu hướng tới Android, chưa được mô tả như một hệ
thống thương mại và chưa tích hợp thanh toán.

Dữ liệu 21 lớp mất cân bằng mạnh: trong validation V3, shirt và tshirt có
hơn 1.500 đối tượng mỗi lớp, trong khi boat và zipper chỉ có hai đối tượng,
glasses có chín đối tượng. Một số nhãn của bộ dữ liệu cũ trộn detection với
segmentation hoặc có toạ độ ngoài miền chuẩn hoá; Ultralytics đã bỏ qua 148
ảnh train và 20 ảnh validation. Ngoài ra, các ánh xạ
\code{Rompers\_Jumpsuits} sang dress và \code{Suiting} sang jacket chỉ là
xấp xỉ ngữ nghĩa.

Đánh giá hiện tại chưa có benchmark thời gian phản hồi trên thiết bị thật.
Sáu test widget Flutter chưa đồng bộ với giao diện hiện hành, và quy tắc bảo
mật Firestore mới được đối chiếu ở mức cấu hình, chưa có biên bản kiểm thử
chéo bằng hai tài khoản độc lập. Đây là các giới hạn của kết quả kiểm thử,
không được tính là chức năng đã xác nhận hoàn toàn.

% ---------------------------------------------------------------------
%  CHƯƠNG 6 — KẾT LUẬN VÀ HƯỚNG PHÁT TRIỂN
% ---------------------------------------------------------------------
\chapter{Kết luận và hướng phát triển}
\label{ch:ketluan}

\begin{muctieu}
Chương này tổng kết phạm vi giải pháp, nêu giới hạn thực nghiệm, rút ra bài
học và đề xuất các hướng phát triển tiếp theo.
\end{muctieu}

\section{Kết luận}

Đồ án ``Xây dựng hệ thống gợi ý phối đồ'' xây dựng ứng dụng Android hỗ trợ
phân tích ảnh trang phục, gợi ý theo hồ sơ và bối cảnh, hỏi đáp, quản lý hồ sơ
cơ thể và thử trang phục theo tư thế camera. Giải pháp gồm Flutter, FastAPI,
YOLO, Gemini, MediaPipe Pose, Open-Meteo và Firebase. Firestore Rules đã dành
đường dẫn tủ đồ riêng cho từng UID, nhưng mã Flutter hiện chưa có luồng sử
dụng đường dẫn này và cũng chưa triển khai danh sách yêu thích; vì vậy hai
chức năng không được tính là kết quả đã hoàn thành.

Đối với thành phần nhận diện, đồ án đã huấn luyện checkpoint YOLO26n theo ba
vòng V1--V3 trên taxonomy 21 lớp. Phiên bản V3 kết hợp dữ liệu Roboflow và
DeepFashion, đạt kết quả validation tốt nhất mAP@0,5:0,95 là 36,69\% trong
lần chạy full; năm epoch fine-tune cuối đạt precision 55,18\%, recall
49,19\%, mAP@0,5 là 49,53\% và mAP@0,5:0,95 là 36,10\%. Các con số này xác
nhận quy trình huấn luyện hoạt động nhưng cũng cho thấy chất lượng chưa đủ để
kết luận đạt mục tiêu 80\% ban đầu.

Backend mặc định dùng \code{backend/models/best.pt}. Kiểm thử tự động ghi
nhận 104/104 test backend đạt; Flutter có 139 test đạt và 6 test widget lỗi.
Do chưa đo trên thiết bị Android thật, báo cáo chưa kết luận về mục tiêu thời
gian phản hồi không quá ba giây.

\section{Bài học rút ra}

Đề tài cho thấy chất lượng của một hệ thống trí tuệ nhân tạo không chỉ phụ
thuộc vào kiến trúc mô hình mà còn phụ thuộc dữ liệu, cách chia tập, quy
trình tiền xử lý và khả năng tích hợp với ứng dụng. Một mô hình có độ chính
xác cao nhưng phản hồi chậm hoặc khó triển khai vẫn có thể không phù hợp với
thiết bị di động.

Việc tách ứng dụng, API, dữ liệu và mô hình thành các thành phần rõ trách
nhiệm giúp kiểm thử và thay thế từng phần dễ hơn. Đồng thời, chức năng thử
đồ 2D minh hoạ giới hạn quan trọng của AI và xử lý ảnh: kết quả trực quan
không đồng nghĩa với mô phỏng chính xác kích cỡ hay đặc tính vật lý của
trang phục.

Quá trình V3 cũng cho thấy dữ liệu cần được kiểm tra trước huấn luyện: một số
nhãn cũ trộn định dạng detection và segmentation, vài hộp vượt miền chuẩn hoá
và phân bố lớp chênh lệch lớn. Việc chia DeepFashion theo \code{item\_id}
giúp tránh rò rỉ các góc chụp của cùng sản phẩm, nhưng ánh xạ giữa taxonomy
DeepFashion và 21 lớp đích vẫn tạo ra nhiễu nhãn.

\section{Hướng phát triển}

Một số hướng phát triển trực tiếp từ các hạn chế ở
Chương~\ref{ch:ketqua} gồm:
\begin{itemize}[leftmargin=1.6em, itemsep=3pt]
  \item mở rộng và cân bằng dữ liệu, bổ sung ảnh thực tế với nhiều nền, tư
        thế và điều kiện ánh sáng;
  \item nhận diện nhiều món đồ trong cùng ảnh và bổ sung thuộc tính như chất
        liệu, hoạ tiết, mùa hoặc hoàn cảnh sử dụng;
  \item cá nhân hoá gợi ý từ phản hồi, lịch sử lựa chọn và phong cách của
        từng người dùng;
  \item thay lớp ảnh 2D bằng phương pháp căn chỉnh theo tư thế hoặc mô hình
        thử đồ sinh ảnh, đồng thời đánh giá rủi ro sai lệch hình ảnh;
  \item hoàn thiện bảo mật, tối ưu thời gian phản hồi, bổ sung iOS và triển
        khai dịch vụ trên hạ tầng có khả năng mở rộng.
\end{itemize}

Mỗi hướng phát triển cần có dữ liệu, tiêu chí đánh giá và kế hoạch kiểm thử
riêng; không nên coi việc thay mô hình mạnh hơn là đủ để giải quyết mọi hạn
chế của hệ thống.


% Chương tra cứu quy cách trình bày. XOÁ DÒNG DƯỚI khi nộp bản chính
% thức, vì đây chỉ là sổ tra cứu chứ không phải một chương của báo cáo.
%% ---------------------------------------------------------------------
%  CHƯƠNG 2 — QUY CÁCH TRÌNH BÀY
%  Chương này đóng vai trò sổ tra cứu: mỗi mục là một mẫu em copy rồi sửa.
%  Khi viết đồ án thật, em thay cả chương này bằng nội dung của mình.
% ---------------------------------------------------------------------
\chapter{Quy cách trình bày trong báo cáo}
\label{ch:soanthao}

\begin{muctieu}
Chương này gom sẵn mẫu cho mọi thành phần em cần dùng: hình, bảng, công
thức, mã nguồn, giải thuật, trích dẫn, chú thích và các hộp nhấn mạnh. Mỗi
mục có đoạn mã mẫu, em copy rồi thay nội dung.
\end{muctieu}

\begin{luuy}
Khi bắt đầu viết đồ án thật, hãy xoá cả chương này. Nó chỉ là sổ tra cứu,
không phải một chương của báo cáo.
\end{luuy}

\section{Chèn hình ảnh}

Đặt file ảnh vào thư mục \code{figs/} rồi dùng mẫu dưới đây. Mọi hình đều
phải có chú thích, có nhãn, và phải được nhắc tới trong thân bài như
Hình~\ref{fig:vidu}.

\begin{figure}[H]
\centering
\includegraphics[width=0.72\textwidth]{dau-campus.png}
\caption{Một ví dụ về chèn hình ảnh từ file.}
\label{fig:vidu}
\end{figure}

Nếu hình là sơ đồ do em tự vẽ thì nên vẽ thẳng bằng TikZ như
Hình~\ref{fig:luongxuly}, vì sơ đồ vẽ bằng TikZ luôn nét khi phóng to và
dùng đúng bộ màu của trường.

\begin{figure}[H]
\centering
% \resizebox ép sơ đồ vừa đúng bề ngang trang, tránh tràn lề
\resizebox{\textwidth}{!}{%
\begin{tikzpicture}[node distance=8mm and 12mm]
  \node[io] (in) {Ảnh đầu vào};
  \node[process, right=of in] (pre) {Tiền xử lý};
  \node[decision, right=of pre] (chk) {Đủ chất lượng?};
  \node[process, right=of chk] (mod) {Mô hình CNN};
  \node[startstop, below=of chk] (loai) {Loại bỏ};
  \draw[flowarr] (in)  -- (pre);
  \draw[flowarr] (pre) -- (chk);
  \draw[flowarr] (chk) -- node[above, font=\scriptsize] {có} (mod);
  \draw[flowarr] (chk) -- node[right, font=\scriptsize] {không} (loai);
\end{tikzpicture}}
\caption{Một ví dụ sơ đồ luồng xử lý vẽ bằng TikZ.}
\label{fig:luongxuly}
\end{figure}

\section{Chèn bảng}

Bảng dùng chú thích đặt phía trên, đúng quy cách trình bày khoa học. Dùng
\code{booktabs} với ba loại đường kẻ \code{toprule}, \code{midrule},
\code{bottomrule}, tránh kẻ dọc cho mọi ô.

\begin{table}[H]
\centering
\caption{Một ví dụ về bảng so sánh kết quả.}
\label{tab:sosanh}
\begin{tabular*}{\textwidth}{@{\extracolsep{\fill}}>{\bfseries\color{daunavy}}l cccc}
\toprule
\rowcolor{gray100}\bfseries\color{daunavy}Mô hình & \bfseries\color{daumaroon}Accuracy & \bfseries\color{daumaroon}Precision & \bfseries\color{daumaroon}Recall & \bfseries\color{daumaroon}F1 \\
\midrule
Huấn luyện từ đầu   & 0,848 & 0,841 & 0,836 & 0,838 \\
Học chuyển giao     & 0,912 & 0,907 & 0,903 & 0,905 \\
\bottomrule
\end{tabular*}
\end{table}

\begin{luuy}
Bảng rộng quá lề là lỗi hay gặp nhất. Nếu bảng tràn, hãy giảm cỡ chữ bằng
\code{small} hoặc \code{footnotesize}, hoặc gộp bớt cột. Đừng để LaTeX tự
tràn rồi bỏ qua.
\end{luuy}

\section{Chèn công thức toán học}

Công thức đặt giữa dòng và có đánh số thì dùng môi trường \code{equation},
như Công thức~\eqref{eq:softmax}. Công thức ngắn chèn trong câu thì viết
giữa hai dấu đô la, ví dụ $y = f(x)$.

\begin{equation}
p(y = k \mid \mathbf{x}) = \frac{\exp(z_k)}{\sum_{j=1}^{K} \exp(z_j)},
\label{eq:softmax}
\end{equation}

trong đó $z_k$ là điểm số thô mà mạng gán cho lớp $k$, và $K$ là số lớp. Mọi
ký hiệu xuất hiện trong công thức đều phải được giải thích ngay sau đó, như
câu này đang làm với \eqref{eq:softmax}.

\begin{ghinho}
Công thức là một phần của câu văn, không phải một khối trôi nổi. Trước công
thức phải có câu dẫn vào, sau công thức phải có câu giải thích ký hiệu hoặc
ý nghĩa. Công thức đứng một mình không ai đọc.
\end{ghinho}

\section{Chèn mã nguồn}

Mẫu này dùng gói \code{listings} thay cho \code{minted} của bản gốc, vì
\code{listings} không cần bật \code{-shell-escape} nên build được ngay trên
Overleaf lẫn máy cá nhân.

Cách thứ nhất là gõ trực tiếp, kết quả như Mã nguồn~\ref{lst:tructiep}.

\begin{lstlisting}[language=Python, caption={Chèn mã nguồn trực tiếp vào báo cáo.}, label={lst:tructiep}]
import torch.nn as nn

# A small convolutional block: conv, batch norm, activation
def conv_block(in_ch, out_ch):
    return nn.Sequential(
        nn.Conv2d(in_ch, out_ch, kernel_size=3, padding=1),
        nn.BatchNorm2d(out_ch),
        nn.ReLU(inplace=True),
    )
\end{lstlisting}

Cách thứ hai là nạp mã từ một file có sẵn trong đồ án. Cách này tiện vì báo
cáo luôn khớp với mã thật, sửa mã là báo cáo tự cập nhật theo. Mã
nguồn~\ref{lst:tufile} được nạp từ file \code{code/vi-du-huan-luyen.py}.

\lstinputlisting[language=Python, caption={Chèn mã nguồn từ file có sẵn.}, label={lst:tufile}]{code/vi-du-huan-luyen.py}

\begin{luuy}
Quy ước của khoa: comment trong code viết bằng tiếng Anh, còn văn xuôi của
báo cáo viết bằng tiếng Việt. Không đặt tiếng Việt bên trong khối mã nguồn,
vì phông mã nguồn không có đủ dấu tiếng Việt.
\end{luuy}

Chỉ đưa vào thân bài những đoạn mã cốt lõi. Mã dài để ở phụ lục hoặc chỉ dẫn
tới kho \code{git} của đồ án.

\section{Biểu diễn giải thuật}

Giải thuật trình bày bằng mã giả, không phải mã của một ngôn ngữ cụ thể,
để người đọc nắm được ý tưởng mà không cần biết ngôn ngữ đó.
Giải thuật~\ref{alg:nguyento} là một ví dụ.

\begin{algorithm}[H]
\caption{Kiểm tra một số tự nhiên có phải số nguyên tố hay không}
\label{alg:nguyento}
\hspace*{\algorithmicindent}\textbf{Đầu vào:} số tự nhiên $n$\\
\hspace*{\algorithmicindent}\textbf{Đầu ra:} \textsc{True} nếu $n$ là số nguyên tố, ngược lại \textsc{False}
\begin{algorithmic}[1]
\Function{LaSoNguyenTo}{$n$}
  \If{$n < 2$}
    \State \Return \textsc{False}
  \EndIf
  \For{$i \gets 2$ \textbf{to} $\lfloor \sqrt{n} \rfloor$}
    \If{$n \bmod i = 0$}
      \State \Return \textsc{False}
    \EndIf
  \EndFor
  \State \Return \textsc{True}
\EndFunction
\end{algorithmic}
\end{algorithm}

\section{Trích dẫn tài liệu tham khảo}

Tài liệu tham khảo khai báo trong file \code{refs.bib}, mỗi tài liệu một
khối. Trong thân bài em gọi bằng lệnh \code{cite} kèm khoá của tài liệu, ví
dụ mô hình AlexNet~\cite{Krizhevsky2012} và mạng ResNet~\cite{He2016}. LaTeX
tự đánh số và tự dựng danh mục ở cuối báo cáo, nên em không phải đánh số tay.

\begin{ghinho}
Chỉ liệt kê trong \code{refs.bib} những tài liệu em thực sự trích dẫn. Tài
liệu khai báo mà không gọi \code{cite} sẽ không xuất hiện trong danh mục.
\end{ghinho}

\section{Chú thích cuối trang}

Khi cần giải thích thêm mà không muốn làm gãy mạch câu, dùng lệnh
\code{footnote}\footnote{Đây là một dòng chú thích cuối trang.}. Đừng lạm
dụng, mỗi trang nhiều lắm một đến hai chú thích.

\section{Các hộp nhấn mạnh}

Mẫu này có sẵn bốn loại hộp để làm nổi những ý quan trọng.

\begin{dinhnghia}[Học chuyển giao]
Học chuyển giao \en{(transfer learning)} là kỹ thuật tận dụng tri thức mà mô
hình đã học được trên một bài toán nguồn có nhiều dữ liệu, để giải một bài
toán đích có ít dữ liệu hơn.
\end{dinhnghia}

\begin{ghinho}
Hộp ghi nhớ dùng cho ý chốt mà em muốn người đọc mang theo sau khi đọc xong
mục đó.
\end{ghinho}

\begin{luuy}
Hộp lưu ý dùng cho cảnh báo, điều kiện áp dụng, hoặc lỗi mà người đọc dễ mắc
nếu làm theo.
\end{luuy}

\begin{vidu}[minh hoạ một tình huống cụ thể]
Hộp ví dụ tự đánh số theo chương, tiện khi em cần dẫn lại một ví dụ ở chỗ
khác trong báo cáo.
\end{vidu}


% ---------------------------------------------------------------------
%  TÀI LIỆU THAM KHẢO VÀ PHỤ LỤC
% ---------------------------------------------------------------------
\backmatter

% \phantomsection tạo mốc neo riêng cho mục lục, nếu thiếu thì hyperref
% báo trùng neo giữa dấu trang và dấu trang cha.
\cleardoublepage
\phantomsection
\addcontentsline{toc}{chapter}{Tài liệu tham khảo}
\bibliographystyle{IEEEtran}
\bibliography{refs}

% ---------------------------------------------------------------------
%  PHỤ LỤC
% ---------------------------------------------------------------------
% Trong \backmatter thì \chapter đã tự thêm vào mục lục, không cần
% \addcontentsline nữa, thêm vào là trùng neo dấu trang.
\chapter{Phụ lục}
\markboth{Phụ lục}{}

\section*{A. Hướng dẫn cài đặt và chạy chương trình}
\phantomsection
\addcontentsline{toc}{section}{A. Hướng dẫn cài đặt và chạy chương trình}

Mã nguồn sử dụng Dart SDK từ 3.13.1, Flutter, Android SDK và Python. Backend
đã được kiểm thử với Python 3.13.7. Quy trình cài đặt được thực hiện từ thư
mục gốc của dự án như sau:

\begin{enumerate}[leftmargin=1.8em, itemsep=4pt]
  \item Tạo môi trường Python bằng lệnh
        \code{python -m venv backend/.venv}. Sau khi kích hoạt môi trường,
        cài các gói trong \path{backend/requirements.txt} bằng \code{pip}.
  \item Tạo \path{backend/.env} từ tệp mẫu cùng thư mục, sau đó cấu hình
        \code{GEMINI\_API\_KEY}. Checkpoint YOLO nằm tại
        \path{backend/models/best.pt}; script tải mô hình pose nằm tại
        \path{backend/scripts/download_pose_model.py}.
  \item Từ thư mục \path{backend}, khởi chạy API bằng
        \code{python -m uvicorn main:app} với host \code{0.0.0.0} và cổng
        \code{8000}. Trạng thái dịch vụ được kiểm tra tại route
        \code{/health}.
  \item Cài phụ thuộc Flutter bằng \code{flutter pub get}. Cấu hình Firebase
        của ứng dụng nằm trong \path{android/app/google-services.json} và
        \path{lib/firebase_options.dart}.
  \item Ứng dụng sử dụng địa chỉ API trong
        \path{lib/config/api_config.dart}; giá trị hiện tại là
        \code{http://192.168.1.25:8000}. Sau khi máy tính và điện thoại kết
        nối cùng mạng, ứng dụng được khởi chạy bằng \code{flutter run}.
\end{enumerate}

\section*{B. Bảng tự kiểm trước khi nộp}
\phantomsection
\addcontentsline{toc}{section}{B. Bảng tự kiểm trước khi nộp}

Các nội dung dưới đây tổng hợp trạng thái hoàn thiện của tài liệu tại thời
điểm biên soạn.

\begin{tukiem}
  \item Cấu trúc chương, mục, trang bìa và quy tắc đánh số của mẫu báo cáo
        được giữ nguyên.
  \item Hình và bảng trong nội dung chính có chú thích, nhãn và tham chiếu
        trong phần thuyết minh liên quan.
  \item Các số liệu YOLO được đối chiếu với tệp kết quả huấn luyện; số liệu
        kiểm thử được ghi theo đầu ra của backend và Flutter.
  \item Những chức năng chưa xuất hiện trong mã nguồn được nêu là giới hạn,
        không được trình bày như kết quả đã hoàn thành.
  \item Tài liệu tham khảo được quản lý bằng \code{refs.bib} và được trích
        dẫn tại nội dung liên quan.
  \item Báo cáo được biên dịch bằng XeLaTeX và kiểm tra lại các tham chiếu
        chéo trước khi tạo bản PDF.
\end{tukiem}

\section*{C. Nhận xét của giảng viên hướng dẫn}
\phantomsection
\addcontentsline{toc}{section}{C. Nhận xét của giảng viên hướng dẫn}

\khungnhanxet[7.5cm]

\vspace{0.8em}
\noindent\begin{minipage}[t]{0.5\textwidth}
\color{daunavy}
{\color{daumaroon}\bfseries Điểm số}

\vspace{0.7em}
\dongchamnhan{Bằng số:}

\vspace{0.9em}
\dongchamnhan{Bằng chữ:}
\end{minipage}\hfill
\okyten{\dauGVHD}{Giảng viên hướng dẫn}


\end{document}
